% Generated from verso/Paper/Text.lean by make -C paper. Do not edit.
\documentclass[11pt]{amsart}

\usepackage[T1]{fontenc}
\usepackage{lmodern}
\usepackage{amsmath,amssymb,amsthm}
\usepackage{mathtools}
\usepackage{booktabs}
\usepackage{array}
\usepackage{enumitem}
\usepackage[margin=1.15in]{geometry}
\usepackage{microtype}
\raggedbottom
\setlength{\emergencystretch}{2.5em}
\setlength{\skip\footins}{8pt plus 4pt minus 2pt}
\usepackage{xcolor}
\usepackage[hidelinks]{hyperref}
\setlist[itemize]{leftmargin=2em}

\newtheorem{theorem}{Theorem}[section]
\newtheorem{proposition}[theorem]{Proposition}
\newtheorem{lemma}[theorem]{Lemma}
\newtheorem{corollary}[theorem]{Corollary}
\theoremstyle{definition}
\newtheorem{definition}[theorem]{Definition}
\newtheorem{remark}[theorem]{Remark}
\newtheorem{example}[theorem]{Example}

\newcommand{\Q}{\mathbb{Q}}
\newcommand{\Z}{\mathbb{Z}}
\newcommand{\C}{\mathbb{C}}
\newcommand{\F}{\mathbb{F}}
\newcommand{\K}{K}
\newcommand{\Kp}{K^{+}}
\newcommand{\cC}{\mathcal{C}}
\newcommand{\Hom}{\operatorname{Hom}}
\newcommand{\End}{\operatorname{End}}
\newcommand{\Tr}{\operatorname{Tr}}
\newcommand{\id}{\mathrm{id}}
\newcommand{\one}{\mathbf 1}
\newcommand{\norm}[1]{\lVert #1\rVert}
\newcommand{\ob}[1]{\texttt{#1}}
\newcommand{\Uq}{U_q(\mathfrak f_4)}
\newcommand{\CF}{\cC(F_4,4)}

\newcommand{\kk}{k}
\newcommand{\Lf}{L}
\newcommand{\Kt}{\Q(\zeta_{13})}
\newcommand{\Fn}{F^{\natural}}
\newcommand{\Rn}{R^{\natural}}

\title[An explicit presentation of $\cC(F_4,4)$]{An explicit presentation of the modular tensor category $\cC(F_4,4)$}
\author{Kim Morrison}
\date{November 2026}

\begin{document}
\begin{abstract}
This paper was written by AIs under loose supervision by Kim Morrison. It probably should only be
read by other AIs. The take-away is that we give explicit F- and R-matrices for $\CF$, and prove that Gannon's
exceptional automorphism of its fusion ring lifts to a braided equivalence between $\CF$ and its
Galois conjugate at $q^5$, whose square is a braid-reversing autoequivalence of $\CF$, settling a
prediction of Morrison--Snyder--Thurston.

$\CF$, the category of the quantum group $\Uq$ at a root of unity of order $52$ (level $4$), has
$16$ simple objects and fusion multiplicities up to $10$; there are $58{,}740$ F-matrices,
with $54{,}881{,}076$ entries, and $2{,}302$ R-matrices. In a normal form that depends only on the
gauge class, the F-symbols lie in $\Q(\zeta_{13})^+$ and the R-symbols in $\Q(\zeta_{13})$, the
smallest fields with the labels of the objects fixed. Allowing the Galois group to permute the objects
by Gannon's automorphism, the braided category has a unique form over the cubic subfield of
$\Q(\zeta_{13})$; it is not split, and three of its simple objects have quaternion division algebras as
endomorphism algebras. We prove exactly, by modular checks with
height bounds, that the data satisfy the hexagon equations and the pentagon equation, and that the twists and the
$S$-matrix computed from them are those of Kac and Peterson, with $S$ invertible. The data and the
programs are available with the paper.
\end{abstract}
\maketitle

\section{Introduction}
A fusion category is described, up to equivalence, by its fusion rules together with its
F-symbols: the matrix coefficients of the associativity constraint in chosen bases of the spaces
$\Hom(X\otimes Y,Z)$ of simple objects. For a braided category one adds R-symbols, the matrix
coefficients of the braiding.

For the categories $\cC(\mathfrak g,k)$ coming from quantum groups at roots of unity (equivalently
from level $k$ integrable representations of affine Lie algebras), explicit F-symbols are known in
closed form only in special cases: for $\mathfrak{sl}_2$ through the $q$-Racah formula, and for
the family $\mathfrak{so}(2p+1)_2$, for all $p$, by Ardonne, Finch and
Titsworth~\cite{ArdonneFinchTitsworth}. Ardonne and Slingerland~\cite{ArdonneSlingerland}
computed $q$-Clebsch--Gordan coefficients and $6j$-symbols for several rank $2$ examples,
including the multiplicity $2$ case $\mathfrak{su}(3)_3/\Z_3$, and Ardonne's Mathematica package
\texttt{alatc}~\cite{Alatc} computes F- and R-symbols from the quantum groups of all untwisted
affine Lie algebras, numerically and in exact form, by this approach. Programs that solve the pentagon and hexagon equations from the fusion
rules (SageMath~\cite{SageFMatrix}, Anyonica~\cite{VercleyenThesis}) are restricted to
multiplicity free fusion rules, and the AnyonWiki~\cite{AnyonWiki} lists the multiplicity free
fusion categories up to rank $7$. Exact F-symbols with multiplicities at a larger scale have been
obtained for Drinfeld centres by M\"aurer and Thiel~\cite{MaurerThiel} and, for the centre of the
Haagerup subfactor, by M\"aurer, Thiel and Vercleyen~\cite{MaurerThielVercleyen}, with
TensorCategories.jl~\cite{TensorCategoriesJl}. VerlindeDB~\cite{VerlindeDB} lists fusion rules and
modular data of many $\cC(\mathfrak g,k)$, including type $F_4$, but no F-symbols. We searched arXiv
titles and abstracts, the forward citations of the papers above, and the databases above, for
F-symbols of $\CF$ or of a category known to be equivalent to it, and found none; we are not aware of
a previous computation of the F-symbols of $\CF$. Galindo~\cite{Galindo} proves that the fusion
category underlying every $\cC(\mathfrak g,k)$ has a split form over the maximal real subfield of the
cyclotomic field of the quantum parameter; we compare this with our fields in
Section~\ref{sec:normalform}.

Gannon found that the fusion ring of $\CF$ has an automorphism $\pi_4$ of order $4$~\cite{Gannon}.
Edie-Michell, who determined the autoequivalence groups of $\cC(\mathfrak g,k)$ for types $A$, $B$, $C$
and $G$, wrote that ``the exceptional auto-equivalence of $\mathcal C(\mathfrak f_4,4)$ remains
completely mysterious''~\cite[Section~1]{EdieMichell}, and Morrison, Snyder and Thurston predicted,
assuming their quantum sufficiency conjecture, a braided equivalence between $\CF$ and a Galois
conjugate realising $\pi_4$, and a braid-reversing autoequivalence~\cite[Section~6]{MorrisonSnyderThurston}.
From the data we prove (Theorem~\ref{thm:autoeq}): $\pi_4$ lifts to a braided equivalence from $\CF$ at
$q$ to $\CF$ at $q^5$, and $\pi_4$ and $\pi_4^3$ do not lift to autoequivalences; $\pi_4^2$ lifts to a
monoidal autoequivalence that reverses the braiding; the monoidal autoequivalence group is $\Z/2$ and
the braided one is trivial; the underlying fusion category has exactly two braidings; and the $24$
Galois conjugates of $\CF$ form three braided equivalence classes. These three classes are the base
changes of one braided fusion category over the cubic subfield $E$ of $\Q(\zeta_{13})$, unique up to
braided equivalence, in which the simple objects over $\ob{0002}$, $\ob{0011}$, $\ob{0101}$ have
quaternion division algebras as endomorphism algebras; the fusion category over $\Q(\zeta_{13})^+$
defined by our F-symbols has no form over $E$ (Section~\ref{sec:forms}).

The category $\CF$ has $16$ simple objects, all self-dual, and its fusion multiplicities reach
$10$ (in $X\otimes X$ for the $4096$-dimensional object $X$ of highest weight $(0,0,1,1)$). Its
F-symbols with all multiplicity indices number about $5.5\times10^7$. We compute them from the
quantum group, as follows.
\begin{itemize}
\item Everything is computed over finite fields $\F_p$ with $p\equiv1\pmod{52}$, which contain
the root of unity. Simple modules of $\Uq$ are built weight space by weight space, and only the
tensor powers $V^{\otimes n}$, $n\le4$, of the $26$-dimensional module $V$ are used, through a
path model in which morphisms act by sparse matrices assembled from a small amount of local data
(the \emph{cells}: $4{,}472$ numbers).
\item Negligible morphisms are removed at the level of the cells, by an explicit description of
the highest weight vectors in $L(A)\otimes V\otimes V$.
\item Every simple object is cut out of a tensor power of $V$ by polynomials in ribbon twists,
and every multiplicity space receives an auxiliary basis, in which the F-symbols with third leg $V$
are computed from the path model.
\item These are brought, modulo $p$, to a normal form defined by an explicit pivot rule, which depends
only on the gauge class and commutes with the Galois action; so it is defined over the field of
moduli, $\Q(\zeta_{13})^+$.
\item All other F-symbols follow from those with third leg $V$ by an explicit pentagon recursion, run
once in the normal-form gauge, which also checks itself; the R-symbols are computed in the same gauge,
over $\Q(\zeta_{13})$.
\item The normal-form data are reconstructed over these fields by Chinese remaindering and rational
reconstruction.
\end{itemize}
The same programs reproduce the closed formulas of Ardonne, Finch and
Titsworth~\cite{ArdonneFinchTitsworth} for $\cC(\mathfrak{so}(13),2)$ (Section~\ref{sec:limits}).

\subsection*{Results}
Let $\zeta=e^{2\pi i/52}$ and $q=\zeta$ (Section~\ref{sec:category}). We use the fields
\[
\K=\Q(\zeta)\supset\Lf=\K^+=\Q(\zeta+\zeta^{-1}),\qquad
\Kt=\Q(\zeta^4)\supset\kk=\Kt^+,
\]
of degrees $24$, $12$, $12$ and $6$. The data consist of the F-symbols $F$ over $\kk$ and the
R-symbols $R$ over $\Kt$ of $\CF$ in the normal-form gauge (Section~\ref{sec:nfrule}).

\begin{theorem}\label{thm:main}
For the fusion rules of $\CF$, the data $(F,R)$ have the following properties.
\begin{enumerate}[label=(\roman*)]
\item Both hexagon equations~\eqref{eq:H1},~\eqref{eq:H2} hold exactly, for all four-tuples of
simple objects and all multiplicity indices.
\item The pentagon equation~\eqref{eq:pentagon} holds exactly for all five-tuples of simple objects
and all multiplicity indices.
\item Every F-matrix and every R-matrix is invertible, every F-matrix with the unit object among
its first three legs is the identity under the natural identification of its labels, and every
object $X$ is self-dual with $F[X,X,X,X]_{(\one,0,0),(\one,0,0)}=1/d_X$, $d_X$ the quantum
dimension~\eqref{eq:qdim}.
\item The numbers
\[
   \tilde\theta_X=d_X^{-1}\sum_Zd_Z\Tr R[X,X,Z],\qquad
   \tilde S_{XY}=\sum_Zd_Z\Tr\bigl(R[Y,X,Z]\,R[X,Y,Z]\bigr)
   \]
are exactly the twists $q^{(X,X+2\rho)}$ and the ratios $S_{XY}/S_{\one\one}$ of the Kac--Peterson
$S$-matrix, and $R[Y,X,Z]\,R[X,Y,Z]=\tilde\theta_Z\tilde\theta_X^{-1}\tilde\theta_Y^{-1}\cdot\id$ for
all $X,Y,Z$. The matrix $\tilde S$ is invertible: $\tilde S^2=\mathcal D^2\cdot\id$ with
$\mathcal D^2=\sum_Xd_X^2$.
\item With the labels of the simple objects fixed, no gauge has all F-symbols in a proper subfield
of $\kk$, and no gauge has all F- and R-symbols in a proper subfield of $\Kt$.
\end{enumerate}


Consequently the data define a braided fusion category over $\Kt$, with F-symbols over $\kk$ and
the fusion rules of $\CF$ (rigid by (iii) and Proposition~\ref{prop:rigid}), and with its spherical
structure of positive dimensions it is a modular tensor category with dimensions $d_X$, twists
$q^{(X,X+2\rho)}$ and the Kac--Peterson $S$-matrix (Proposition~\ref{prop:modular}).
\end{theorem}

Theorem~\ref{thm:main} is a statement about the data files. That the category they define is $\CF$
rests on the construction, as discussed in Section~\ref{sec:limits}.

\begin{table}[ht]
\centering
\small
\begin{tabular}{lr}
\toprule
field of F; of R & $\kk$ (degree 6); $\Kt$ (degree 12) \\
F-matrices; entries (including zeros) & $58{,}740$; $54{,}881{,}076$ \\
largest F-matrix; largest multiplicity & $465\times465$; $10$ \\
nonzero rational coefficients of F (basis $\eta_1,\dots,\eta_6$) & $328{,}727{,}941$ \\
largest numerator or denominator, F & $99$ bits \\
R-matrices; entries & $2{,}302$; $8{,}116$ \\
largest numerator or denominator, R; $R^{-1}$ & $47$ bits; $41$ bits \\
five-tuples with $\Hom(A\otimes B\otimes C\otimes D,E)\ne0$ & $1{,}030{,}051$ \\
\bottomrule
\end{tabular}
\medskip
\caption{Sizes of the data.}
\label{tab:sizes}
\end{table}

Statements (i), (ii) and (iv) are proved by checking the identities modulo primes
$p\equiv1\pmod{52}$ at all embeddings into $\F_p$ and applying an explicit height bound
(Lemma~\ref{lem:height}). For the pentagon we use Lemma~\ref{lem:reduction}, which reduces it to the
$63{,}841$ five-tuples whose fourth object is the $26$-dimensional object $V$; larger exact subsets give
additional checks (Section~\ref{sec:pentagon}). Statement (v) rests on the normal form (Section~\ref{sec:fields}).

We also show (Proposition~\ref{prop:integrality}) that no gauge makes all F-symbols of $\CF$
integral at the primes above $13$: the object $(0,0,1,1)$ has a quantum dimension of
valuation $2/3$ there, and a gauge invariant combination of F-symbols equals its inverse.

\subsection*{Outline}
Section~\ref{sec:category} describes $\CF$: root data, the root of unity, the simple objects and the
fusion graph of $V$. Section~\ref{sec:conventions} fixes conventions for F- and R-symbols with
multiplicities, writes out the pentagon and hexagon equations, and proves the reduction of the
pentagon equation to one generator. Section~\ref{sec:construction} describes the construction,
including the normal form, in enough detail to reimplement it. Section~\ref{sec:fields} discusses the
fields and proves Theorem~\ref{thm:main}(v), Section~\ref{sec:reconstruction} the exact reconstruction
and Section~\ref{sec:verification} the verification and unitarity. Section~\ref{sec:autoeq}
determines the autoequivalences and braidings and the forms over the cubic subfield, and Section~\ref{sec:integrality} proves the
non-integrality at $13$. Sections~\ref{sec:data} and~\ref{sec:availability} describe the data files
and how to reproduce the computation, and Section~\ref{sec:limits} discusses limitations and an open
question.

\section{The category}\label{sec:category}
\subsection{Root data}
We use Bourbaki's numbering of the simple roots of $F_4$~\cite{Bourbaki}: $\alpha_1,\alpha_2$
long and $\alpha_3,\alpha_4$ short, with Cartan matrix
\[
\bigl(\langle\alpha_j,\alpha_i^\vee\rangle\bigr)_{i,j}=
\begin{pmatrix}2&-1&0&0\\-1&2&-1&0\\0&-2&2&-1\\0&0&-1&2\end{pmatrix}.
\]
Weights are written by their Dynkin labels $\lambda=(\lambda_1,\lambda_2,\lambda_3,\lambda_4)$,
$\lambda_i=\langle\lambda,\alpha_i^\vee\rangle$, and in the data files as strings such as
\ob{0011}. The invariant form is normalised so that \emph{short} roots have length $2$; then long
roots have length $4$, and $(\lambda,\mu)\in\Z$ for all integral weights in the cases we need. The
Weyl vector is $\rho=(1,1,1,1)$, the highest root is $\theta=(1,0,0,0)$, the comarks are
$(2,3,2,1)$, and the dual Coxeter number is $h^\vee=9$. There are $24$ positive roots.

\subsection{The root of unity}
Let $k=4$, so $k+h^\vee=13$, and let $q$ be a primitive root of unity of order
$\ell=2\cdot2\cdot13=52$ (twice the lacing number times $k+h^\vee$). Over $\C$ we take
$q=\zeta=e^{2\pi i/52}$; the other primitive $52$nd roots give the Galois conjugate categories
(Section~\ref{sec:fields}). With our normalisation of the form, $K_i$ acts on a weight space of
weight $\mu$ by $q^{(\alpha_i,\mu)}$, so $q_i=q^{(\alpha_i,\alpha_i)/2}$ is $q$ for short and
$q^2$ for long simple roots. In the conventions of Rowell's survey~\cite{Rowell}, where the form is normalised in the same way and
$q^2$ is a primitive $\ell'$-th root of unity, this is the uniform case $\ell'=26$, $m=2\mid\ell'$, at
$q=e^{\pi i/26}$ (Section~\ref{sec:unitarity}).

\subsection{The category}
Let $\cC$ be the category of tilting modules of $\Uq$ at $q$ (with divided powers), modulo the
tensor ideal generated by the indecomposable tilting modules of quantum dimension zero. By
Andersen~\cite{Andersen92} and Andersen and Paradowski~\cite{AndersenParadowski} (see the account
in~\cite[Section~3]{Rowell}, which we follow), $\cC$ is a semisimple ribbon category whose simple
objects are labelled by the dominant weights $\lambda$ in the open fundamental alcove
\[
\{\lambda\ \text{dominant}:(\lambda+\rho,\theta)<26\}
=\{\lambda:2\lambda_1+3\lambda_2+2\lambda_3+\lambda_4\le4\},
\]
that is, the integrable highest weights of level at most $4$; by Kazhdan and
Lusztig~\cite{KazhdanLusztigI,KazhdanLusztigII,KazhdanLusztigIII,KazhdanLusztigIV} and
Finkelberg~\cite{Finkelberg} it is equivalent to the category of level $4$ integrable modules of
the affine Lie algebra of type $F_4^{(1)}$. We write $\CF$ for it. There are $16$ simple objects.
Their quantum dimensions and twists are
\begin{equation}\label{eq:qdim}
d_\lambda=\prod_{\alpha>0}\frac{[(\lambda+\rho,\alpha)]}{[(\rho,\alpha)]},\qquad
[n]=\frac{q^n-q^{-n}}{q-q^{-1}},\qquad
\theta_\lambda=q^{(\lambda,\lambda+2\rho)}=e^{2\pi ih_\lambda},
\end{equation}
with conformal weights $h_\lambda=(\lambda,\lambda+2\rho)/52$. All $(\lambda,\lambda+2\rho)$ are even,
so the twists lie in $\mu_{26}$. The central charge is
$c=k\dim\mathfrak f_4/(k+h^\vee)=4\cdot52/13=16$, and the global dimension is
$\mathcal D^2=\sum_\lambda d_\lambda^2\approx7402.7940133506$, $\mathcal D\approx86.0394910105$.
Table~\ref{tab:objects} lists the simple objects (\texttt{exact\_checks.py}).

\begin{table}[ht]
\centering
\renewcommand{\arraystretch}{1.05}
\begin{tabular}{rlrrlrl}
\toprule
 & $\lambda$ & $\dim L(\lambda)$ & $d_\lambda$ & $h_\lambda$ & depth & $v_{13}(d_\lambda)$ \\
\midrule
0 & \ob{0000} & 1 & 1 & 0 & 0 & 0 \\
1 & \ob{0001} & 26 & 7.3021868179 & 6/13 & 1 & 0 \\
2 & \ob{0002} & 324 & 18.3587793437 & 1 & 2 & 0 \\
3 & \ob{0010} & 273 & 19.3587793437 & 12/13 & 2 & 0 \\
4 & \ob{1000} & 52 & 7.3021868179 & 9/13 & 2 & 0 \\
5 & \ob{0003} & 2652 & 19.3587793437 & 21/13 & 3 & 0 \\
6 & \ob{0011} & 4096 & 43.0197455053 & 3/2 & 3 & 2/3 \\
7 & \ob{0100} & 1274 & 19.3587793437 & 18/13 & 3 & 0 \\
8 & \ob{1001} & 1053 & 26.6609661616 & 16/13 & 3 & 0 \\
9 & \ob{0004} & 16302 & 7.3021868179 & 30/13 & 4 & 0 \\
10 & \ob{0012} & 34749 & 26.6609661616 & 28/13 & 4 & 0 \\
11 & \ob{0020} & 19448 & 19.3587793437 & 27/13 & 4 & 0 \\
12 & \ob{0101} & 19278 & 25.6609661616 & 2 & 4 & 0 \\
13 & \ob{1002} & 10829 & 26.6609661616 & 24/13 & 4 & 0 \\
14 & \ob{1010} & 8424 & 26.6609661616 & 23/13 & 4 & 0 \\
15 & \ob{2000} & 1053 & 7.3021868179 & 20/13 & 4 & 0 \\
\bottomrule
\end{tabular}
\medskip
\caption{The simple objects of $\CF$, in the order used in the data files. ``Depth'' is the least
$n$ with $L(\lambda)$ a summand of $V^{\otimes n}$, $V=L(\ob{0001})$. The last column is the
valuation of $d_\lambda$ at either prime of $\K$ above $13$, normalised by $v(13)=1$
(Section~\ref{sec:integrality}).}
\label{tab:objects}
\end{table}

\subsection{Fusion rules and the generator}
The fusion multiplicities $N_{XY}^Z=\dim\Hom(X\otimes Y,Z)$ are given by the Kac--Walton
formula~\cite{Walton90a,Walton90b} (see also~\cite{KacBook}): the classical tensor product
multiplicities, folded by the affine Weyl group action at level $4$ (shifted by $\rho$). We
computed them this way, and checked that they agree with the Verlinde formula applied to the
Kac--Peterson $S$-matrix~\cite{KacPeterson} and that they form an associative commutative ring
with the $d_\lambda$ as a character. Every simple object is self-dual. There are $2{,}302$ triples
with $N_{XY}^Z\ne0$, and the largest multiplicity is
$N_{\ob{0011},\ob{0011}}^{\ob{0011}}=10$.

Let $V=L(\ob{0001})$ be the $26$-dimensional module. Table~\ref{tab:graph} gives $X\otimes V$ for
every simple $X$, that is, the fusion graph of $V$. All multiplicities are $1$ except the two
loops at \ob{0011} and \ob{0012}, of multiplicity $2$. Every simple object occurs in $V^{\otimes4}$
(the depths are in Table~\ref{tab:objects}). In particular
\[
V\otimes V=\one\oplus V\oplus L(\ob{0002})\oplus L(\ob{0010})\oplus L(\ob{1000}),
\]
of dimensions $676=1+26+324+273+52$, where the first three summands lie in $S^2V$ and the last
two in $\Lambda^2V$.

\begin{table}[ht]
\centering
\small
\begin{tabular}{ll}
\toprule
$X$ & $X\otimes V$ \\
\midrule
\ob{0000} & \ob{0001} \\
\ob{0001} & \ob{0000} + \ob{0001} + \ob{0002} + \ob{0010} + \ob{1000} \\
\ob{0002} & \ob{0001} + \ob{0002} + \ob{0010} + \ob{0003} + \ob{0011} + \ob{1001} \\
\ob{0010} & \ob{0001} + \ob{0002} + \ob{0010} + \ob{1000} + \ob{0011} + \ob{0100} + \ob{1001} \\
\ob{1000} & \ob{0001} + \ob{0010} + \ob{1001} \\
\ob{0003} & \ob{0002} + \ob{0003} + \ob{0011} + \ob{0004} + \ob{0012} + \ob{1002} \\
\ob{0011} & \ob{0002} + \ob{0010} + \ob{0003} + 2\,\ob{0011} + \ob{0100} + \ob{1001} + \ob{0012} + \ob{0020} + \ob{0101} + \ob{1002} + \ob{1010} \\
\ob{0100} & \ob{0010} + \ob{0011} + \ob{1001} + \ob{0101} + \ob{1010} \\
\ob{1001} & \ob{0002} + \ob{0010} + \ob{1000} + \ob{0011} + \ob{0100} + \ob{1001} + \ob{1002} + \ob{1010} + \ob{2000} \\
\ob{0004} & \ob{0003} + \ob{0004} + \ob{0012} \\
\ob{0012} & \ob{0003} + \ob{0011} + \ob{0004} + 2\,\ob{0012} + \ob{0020} + \ob{0101} + \ob{1002} \\
\ob{0020} & \ob{0011} + \ob{0012} + \ob{0020} + \ob{0101} + \ob{1010} \\
\ob{0101} & \ob{0011} + \ob{0100} + \ob{0012} + \ob{0020} + \ob{0101} + \ob{1002} + \ob{1010} \\
\ob{1002} & \ob{0003} + \ob{0011} + \ob{1001} + \ob{0012} + \ob{0101} + \ob{1002} + \ob{1010} \\
\ob{1010} & \ob{0011} + \ob{0100} + \ob{1001} + \ob{0020} + \ob{0101} + \ob{1002} + \ob{1010} + \ob{2000} \\
\ob{2000} & \ob{1001} + \ob{1010} \\
\bottomrule
\end{tabular}
\medskip
\caption{The fusion graph of $V=\ob{0001}$: the decomposition of $X\otimes V$ in $\CF$.}
\label{tab:graph}
\end{table}

\section{Conventions for F- and R-symbols}\label{sec:conventions}
We fix conventions for F-symbols with multiplicities; see~\cite{EGNO,BakalovKirillov} for
generalities on fusion and braided categories. Let $\cC$ be a semisimple $k$-linear monoidal
category over a field $k$ of characteristic zero, with simple unit $\one$, finitely many simple
objects up to isomorphism and $\End(X)=k$ for simple $X$. Fix representatives
$\mathrm{Irr}(\cC)$, and for simple $X,Y,Z$ a basis $\phi^{XY\to Z}_i$, $0\le i<N_{XY}^Z$, of
$\Hom(X\otimes Y,Z)$ (the \emph{fusion maps}). We suppress the associator of the ambient category in the
notation (in general one inserts $\alpha_{X,Y,W}$ on the right hand side of~\eqref{eq:F}); our
computations take place in a category of matrices over path spaces in which composition and tensor
product are strictly associative. The two families in~\eqref{eq:F} below are bases of
$\Hom(X\otimes Y\otimes W,T)$, by semisimplicity.

\begin{definition}\label{def:F}
For simple $X,Y,W,T$, the \emph{F-matrix} $F[X,Y,W,T]$ is defined by
\begin{equation}\label{eq:F}
\phi^{ZW\to T}_j\circ(\phi^{XY\to Z}_i\otimes\id_W)
=\sum_{Z',k,l} F[X,Y,W,T]_{(Z,i,j),(Z',k,l)}\;
\phi^{XZ'\to T}_l\circ(\id_X\otimes\phi^{YW\to Z'}_k).
\end{equation}
Rows are labelled by triples $(Z,i,j)$ with $i<N_{XY}^Z$, $j<N_{ZW}^T$, and columns by triples
$(Z',k,l)$ with $k<N_{YW}^{Z'}$, $l<N_{XZ'}^T$.
\end{definition}

\begin{proposition}\label{prop:pentagon}
For all simple $A,B,C,D,E$, all $P,Q,R,S$ and all indices,
\begin{multline}\label{eq:pentagon}
\sum_{l} F[P,C,D,E]_{(Q,i_2,i_3),(R,k,l)}\,F[A,B,R,E]_{(P,i_1,l),(S,m,n)} \\
=\sum_{T,a,b,c} F[A,B,C,Q]_{(P,i_1,i_2),(T,a,b)}\,F[A,T,D,E]_{(Q,b,i_3),(S,c,n)}\,
F[B,C,D,S]_{(T,a,c),(R,k,m)}.
\end{multline}
\end{proposition}

\begin{proof}
Both sides are the coefficient of the basis vector
$\phi^{AS\to E}_n\circ(\id_A\otimes\phi^{BR\to S}_m)\circ(\id_{A\otimes B}\otimes\phi^{CD\to R}_k)$
of $\Hom(A\otimes B\otimes C\otimes D,E)$ in the expansion of
$v=\phi^{QD\to E}_{i_3}\circ(\phi^{PC\to Q}_{i_2}\otimes\id_D)\circ(\phi^{AB\to P}_{i_1}\otimes\id_{C\otimes D})$.
For the left hand side, apply~\eqref{eq:F} to the first two factors of $v$ (with
$(X,Y,W,T)=(P,C,D,E)$), use the interchange law to move $\phi^{AB\to P}_{i_1}$ past
$\phi^{CD\to R}_k$, and apply~\eqref{eq:F} with $(A,B,R,E)$. For the right hand side, apply~\eqref{eq:F}
with $(A,B,C,Q)$, then with $(A,T,D,E)$, and finally, inside the second tensor factor, with
$(B,C,D,S)$.
\end{proof}

In matrix form, for fixed $A,B,C,D,E$ consider the five label sets
\begin{align*}
I_0&=\{(P,i_1,Q,i_2,i_3)\},& I_1&=\{(P,i_1,R,k,l)\},& I_2&=\{(R,k,S,m,n)\},\\
J_1&=\{(T,a,b,Q,i_3)\},& J_2&=\{(T,a,S,c,n)\},
\end{align*}
indexing the bases of $\Hom(A\otimes B\otimes C\otimes D,E)$ for the bracketings
$((AB)C)D$, $(AB)(CD)$, $A(B(CD))$, $(A(BC))D$ and $A((BC)D)$, all of the same size
$n_{ABCDE}=\dim\Hom(A\otimes B\otimes C\otimes D,E)$. Then~\eqref{eq:pentagon} says
\begin{equation}\label{eq:pentagon-matrix}
M_1M_2=N_1N_2N_3,
\end{equation}
where $M_1=\bigoplus_{P,i_1}F[P,C,D,E]$ ($I_0\to I_1$), $M_2=\bigoplus_{R,k}F[A,B,R,E]$
($I_1\to I_2$), $N_1=\bigoplus_{Q,i_3}F[A,B,C,Q]$ ($I_0\to J_1$), $N_2=\bigoplus_{T,a}F[A,T,D,E]$
($J_1\to J_2$) and $N_3=\bigoplus_{S,n}F[B,C,D,S]$ ($J_2\to I_2$), each a block diagonal matrix
with the identity on the multiplicity indices that are not acted on. We call
$(A,B,C,D,E)$ a \emph{five-tuple}; $D$ is its fourth object. There are $1{,}030{,}051$ five-tuples
with $n_{ABCDE}\ne0$ for $\CF$, and $n_{ABCDE}$ reaches $19{,}900$.

\subsection*{Gauge}
A change of fusion maps $\tilde\phi^{XY\to Z}_i=\sum_ag^{XY}_Z[i,a]\phi^{XY\to Z}_a$, with
$g^{XY}_Z$ invertible, changes $F[X,Y,W,T]$ to $G_LF[X,Y,W,T]G_R^{-1}$, with $G_L$ and $G_R$
block diagonal matrices assembled from tensor products of the $g$'s on the two label sets, and
preserves the pentagon equation. In particular, if all four multiplicity spaces involved are one
dimensional,
\begin{equation}\label{eq:gauge-scalar}
\tilde F[X,Y,W,T]_{(Z,0,0),(Z',0,0)}=\frac{g^{XY}_Zg^{ZW}_T}{g^{YW}_{Z'}g^{XZ'}_T}\,
F[X,Y,W,T]_{(Z,0,0),(Z',0,0)} .
\end{equation}

\subsection*{R-symbols and hexagons}
If $\cC$ is braided with braiding $c_{X,Y}:X\otimes Y\to Y\otimes X$, the R-matrices are
defined by
\begin{equation}\label{eq:R}
\phi^{YX\to Z}_j\circ c_{X,Y}=\sum_iR[X,Y,Z]_{j,i}\,\phi^{XY\to Z}_i ,
\end{equation}
square matrices of size $N_{XY}^Z=N_{YX}^Z$. Write $\bar F$ for the inverse of an F-matrix. The two
hexagon axioms, for $c_{X,Y\otimes W}$ and for $c_{X\otimes Y,W}$, become the following
identities, for all simple $X,Y,W,T$ and all indices (for the first, all $Z,k,l$ and
$Z'',k'',l''$; for the second, all $Z,i,l$ and $V',s',t$):
\begin{multline}\label{eq:H1}
\sum_{\substack{U,m,m',n\\V',s,s',t}}F[Y,W,X,T]_{(Z,k,l),(U,m,n)}\,R[X,W,U]_{m,m'}\,
\bar F[Y,X,W,T]_{(U,m',n),(V',s,t)}\\
\cdot R[X,Y,V']_{s,s'}\,F[X,Y,W,T]_{(V',s',t),(Z'',k'',l'')}
=\delta_{ZZ''}\delta_{kk''}\,R[X,Z,T]_{l,l''},
\end{multline}
\begin{multline}\label{eq:H2}
\sum_jR[Z,W,T]_{l,j}\,F[X,Y,W,T]_{(Z,i,j),(V',s',t)}\\
=\sum_{U,m,m',n,s}\bar F[W,X,Y,T]_{(Z,i,l),(U,m,n)}\,R[X,W,U]_{m,m'}\\
\cdot F[X,W,Y,T]_{(U,m',n),(V',s,t)}\,R[Y,W,V']_{s,s'}.
\end{multline}

\begin{proof}[Derivation]
For~\eqref{eq:H2}, apply both sides of $c_{X\otimes Y,W}=(c_{X,W}\otimes\id_Y)\circ(\id_X\otimes
c_{Y,W})$ after the basis vector $\phi^{WZ\to T}_l\circ(\id_W\otimes\phi^{XY\to Z}_i)$ of
$\Hom(W\otimes X\otimes Y,T)$. By naturality, the left side gives
$\phi^{WZ\to T}_l\circ c_{Z,W}\circ(\phi^{XY\to Z}_i\otimes\id_W)
=\sum_jR[Z,W,T]_{l,j}\,\phi^{ZW\to T}_j\circ(\phi^{XY\to Z}_i\otimes\id_W)$, which we expand
by~\eqref{eq:F}. For the right side, rewrite the basis vector in the left-bracketed basis with
$\bar F[W,X,Y,T]$, compose with $c_{X,W}\otimes\id_Y$ using~\eqref{eq:R}, expand with
$F[X,W,Y,T]$, and compose with $\id_X\otimes c_{Y,W}$ using~\eqref{eq:R} again. Comparing
coefficients of $\phi^{XV'\to T}_t\circ(\id_X\otimes\phi^{YW\to V'}_{s'})$ gives~\eqref{eq:H2}.
For~\eqref{eq:H1}, apply $c_{X,Y\otimes W}=(\id_Y\otimes c_{X,W})\circ(c_{X,Y}\otimes\id_W)$
after $\phi^{ZX\to T}_l\circ(\phi^{YW\to Z}_k\otimes\id_X)$; naturality gives
$\sum_{l''}R[X,Z,T]_{l,l''}\,\phi^{XZ\to T}_{l''}\circ(\id_X\otimes\phi^{YW\to Z}_k)$ on one side,
and the other side is expanded with $F[Y,W,X,T]$, $R[X,W,\cdot]$, $\bar F[Y,X,W,T]$,
$R[X,Y,\cdot]$ and $F[X,Y,W,T]$ in turn.
\end{proof}

For the verification we use the equivalent forms without $\bar F$, in which the R-blocks on one
side are replaced by their inverses: with $P_1,P_2,P_3$ the block matrices built from the blocks
$R[X,Y,\cdot]^{-1}$, $R[X,W,\cdot]^{-1}$, $R[X,\cdot,T]^{-1}$ acting on the first, first and last
multiplicity index respectively, and $Q_4,Q_5,Q_6$ those built from $R[Y,W,\cdot]$,
$R[X,W,\cdot]$, $R[\cdot,W,T]$,
\begin{equation}\label{eq:hex-matrix}
\begin{gathered}
P_1\,F[Y,X,W,T]\,P_2=F[X,Y,W,T]\,P_3\,F[Y,W,X,T],\\
Q_5\,F[X,W,Y,T]\,Q_4=F[W,X,Y,T]\,Q_6\,F[X,Y,W,T].
\end{gathered}
\end{equation}
The first is~\eqref{eq:H1} multiplied out and inverted, the second is~\eqref{eq:H2} multiplied on
the left by $F[W,X,Y,T]$. They are equivalent to~\eqref{eq:H1},~\eqref{eq:H2} given that the
F-matrices are invertible and that the inverse R-matrices are the inverses of the R-matrices,
both of which we check.

\subsection*{Units, duals and dimensions}
In our data every F-matrix with the unit object among its first three legs is the identity under
the natural identification of labels: for $F[X,\one,W,T]$ the row $(X,0,j)$ corresponds to the
column $(W,0,j)$, for $F[X,Y,\one,T]$ the row $(T,i,0)$ to the column $(Y,0,i)$, and for
$F[\one,Y,W,T]$ the row $(Y,0,j)$ to the column $(T,j,0)$. For a self-dual simple $X$ with
$N_{XX}^\one=1$ the entry $F[X,X,X,X]_{(\one,0,0),(\one,0,0)}$ transforms
by~\eqref{eq:gauge-scalar} with the factor $g^{\one X}_X/g^{X\one}_X$, so it is invariant under
gauges that keep the unit normalisation; in a spherical category it equals $\nu_X/d_X$, with
$\nu_X$ the Frobenius--Schur indicator. In our data it is exactly $1/d_X$ for all $16$ objects
(Section~\ref{sec:verification}), so $\nu_X=1$ for all $X$, and by the following criterion the
data are rigid.

\begin{proposition}\label{prop:rigid}
Let $\cC$ be as above, with the unit maps $\phi^{\one X\to X}_0$ and $\phi^{X\one\to X}_0$ the
canonical isomorphisms. Suppose that for every simple $X$ there is a simple $X^*$ with
$N_{XX^*}^{\one}=N_{X^*X}^{\one}=1$ and $F[X,X^*,X,X]_{(\one,0,0),(\one,0,0)}\ne0$. Then every
simple object has a left and a right dual, so $\cC$ is rigid.
\end{proposition}

\begin{proof}
Put $a=\phi^{XX^*\to\one}_0$, $b=\phi^{X^*X\to\one}_0$, and let $c$ span
$\Hom(\one,X^*\otimes X)$, which is one dimensional by semisimplicity; then $b\circ c=\beta$ is a
nonzero scalar. By~\eqref{eq:F} with $(X,Y,W,T)=(X,X^*,X,X)$,
\[
a\otimes\id_X=\sum_{Z',k,l}F[X,X^*,X,X]_{(\one,0,0),(Z',k,l)}\,
\phi^{XZ'\to X}_l\circ(\id_X\otimes\phi^{X^*X\to Z'}_k).
\]
Precompose with $\id_X\otimes c$. Since $\phi^{X^*X\to Z'}_k\circ c\in\Hom(\one,Z')$ vanishes unless
$Z'=\one$, the composite $(a\otimes\id_X)\circ(\id_X\otimes c)\in\End(X)$ equals
$s=\beta\,F[X,X^*,X,X]_{(\one,0,0),(\one,0,0)}\ne0$. Replacing $c$ by $c/s$ we get
$(a\otimes\id_X)\circ(\id_X\otimes c)=\id_X$. The other composite
$u=(\id_{X^*}\otimes a)\circ(c\otimes\id_{X^*})\in\End(X^*)$ is a scalar $t$, and
$t\,a=a\circ(\id_X\otimes u)=a\circ\bigl(((a\otimes\id_X)\circ(\id_X\otimes c))\otimes\id_{X^*}\bigr)=a$
by the interchange law, so $t=1$. Thus $(X^*,a,c)$ exhibits $X^*$ as a right dual of $X$. The same
argument for $X^*$, whose partner is $X$, exhibits $X^*$ as a left dual of $X$. Duals of direct sums
are direct sums of duals.
\end{proof}

\subsection{Reducing the pentagon equation to one generator}\label{sec:reduction}
The pentagon equations are highly redundant. The following reduction to the five-tuples whose
fourth object is a fixed generator is well known in spirit: it is the standard argument with the
associahedron from the proof of coherence~\cite{MacLane}. We include a short proof for
completeness, since we could not locate a reference in exactly this form, with multiplicities and
without a unit normalisation.

\begin{lemma}\label{lem:reduction}
Let $I$ be a finite set of labels with nonnegative integers $N_{ab}^c$ ($a,b,c\in I$) satisfying
$\sum_zN_{ab}^zN_{zc}^t=\sum_zN_{bc}^zN_{az}^t$ for all $a,b,c,t$, and let $F[a,b,c,t]$ be
invertible matrices over a field, with rows and columns indexed as in Definition~\ref{def:F}.
Fix $v\in I$ and let $I_v\subseteq I$ be the set of labels $d$ for which there is a path
$v=x_1\to x_2\to\dots\to x_n=d$, $n\ge1$, with $N_{x_jv}^{x_{j+1}}\ne0$ for all $j$ (in a fusion
category: the simple summands of $v^{\otimes n}$, $n\ge1$). If the pentagon
equation~\eqref{eq:pentagon-matrix} holds for all five-tuples $(a,b,c,v,e)$ with $a,b,c,e\in I$,
then it holds for all five-tuples $(a,b,c,d,e)$ with $d\in I_v$. In particular, if $I_v=I$, it
holds for all five-tuples.
\end{lemma}

No unit object, unit normalisation or triangle axiom is involved. If $I$ is the set of simple
objects of a fusion category in which every simple object is a summand of some $v^{\otimes n}$ and
$v$ is self-dual (so that $\one$ is a summand of $v\otimes v$), then $I_v=I$. If the unit is only a
summand of $v^{\otimes0}$, the remaining case $d=\one$ holds as soon as the F-matrices with unit
third leg are identities, since the pentagon for $(a,b,c,\one,e)$ then reads $F[a,b,c,e]=F[a,b,c,e]$.

\begin{proof}
\emph{A model category.} Let $\mathcal M$ be the category of $I$-graded finite-dimensional vector
spaces $X=(X_a)_{a\in I}$, with graded linear maps as morphisms, and with the bifunctor
$(X\otimes Y)_c=\bigoplus_{a,b}X_a\otimes Y_b\otimes H^{ab}_c$, where $H^{ab}_c$ has basis the
multiplicity indices $0,\dots,N_{ab}^c-1$. A label $a$ is identified with the object that is the
ground field in degree $a$ and zero elsewhere. For labels $a,b,c$ the $t$-components of
$(a\otimes b)\otimes c$ and $a\otimes(b\otimes c)$ have bases indexed by the triples $(z,i,j)$ and
$(z',k,l)$ of Definition~\ref{def:F}, and we let $\alpha_{a,b,c}$ act on the $t$-component by the
matrix $F[a,b,c,t]$, on coefficient row vectors, so that composites correspond to matrix products in
the order written. For arbitrary objects put
$\alpha_{X,Y,W}=\sum_{a,b,c}\id_{X_a\otimes Y_b\otimes W_c}\otimes\alpha_{a,b,c}$. This is a
natural isomorphism $(X\otimes Y)\otimes W\to X\otimes(Y\otimes W)$: natural because it acts as the
identity on the spaces $X_a,Y_b,W_c$, and invertible because every $F[a,b,c,t]$ is; no pentagon is
needed for this. For objects $X,Y,W,U$ let $\mathrm P(X,Y,W,U)$ be the statement that the pentagon
diagram for $\alpha$ with these four objects commutes.

\emph{(1) Dictionary.} For labels $a,b,c,d$, on the $e$-component the five maps of the pentagon
diagram, $\alpha_{a\otimes b,c,d}$, $\alpha_{a,b,c\otimes d}$, $\alpha_{a,b,c}\otimes\id_d$,
$\alpha_{a,b\otimes c,d}$ and $\id_a\otimes\alpha_{b,c,d}$, have the matrices $M_1$, $M_2$, $N_1$,
$N_2$, $N_3$ of~\eqref{eq:pentagon-matrix}: the direct sums over $P,R,Q,T,S$ there are exactly the
decompositions of $a\otimes b$, $c\otimes d$, $b\otimes c$, \dots\ into labels, with the identity
on the multiplicity spaces. So $\mathrm P(a,b,c,d)$ holds if and only if~\eqref{eq:pentagon-matrix}
holds for $(a,b,c,d,e)$ for every $e$.

\emph{(2) Additivity.} For arbitrary objects, each of the five maps in $\mathrm P(X,Y,W,U)$ is,
by definition, the sum over labels $a,b,c,d$ of the identity of $X_a\otimes Y_b\otimes W_c\otimes U_d$
tensored with the corresponding map for $(a,b,c,d)$. So $\mathrm P(X,Y,W,U)$ holds if and only if
$\mathrm P(a,b,c,d)$ holds for all labels with $X_a,Y_b,W_c,U_d\ne0$. Also, if $\mathrm P(a,b,c,d)$
holds then so do its tensor products with an identity on the left or on the right, by
functoriality of $\otimes$.

\emph{(3) The associahedron.} For objects $X_1,\dots,X_5$ the bracketings of
$X_1\otimes\dots\otimes X_5$ and the single applications of $\alpha$ to a subterm form the
vertices ($14$) and edges ($21$) of the associahedron $K_5$~\cite{Stasheff,MacLane}. Its nine faces
are six pentagons, $\mathrm P(X_1,X_2,X_3,X_4)\otimes\id_{X_5}$, $\id_{X_1}\otimes\mathrm P(X_2,X_3,X_4,X_5)$,
$\mathrm P(X_1\otimes X_2,X_3,X_4,X_5)$, $\mathrm P(X_1,X_2\otimes X_3,X_4,X_5)$,
$\mathrm P(X_1,X_2,X_3\otimes X_4,X_5)$ and $\mathrm P(X_1,X_2,X_3,X_4\otimes X_5)$, and three
squares. Two squares commute by naturality of $\alpha$ with respect to the morphism $\alpha$ in one of
its arguments, and the third, in which two associators act on disjoint parts of the bracketing,
commutes by functoriality of $\otimes$. The boundary of $K_5$ is a $2$-sphere, so the complement of
one open face is a disc subdivided by the other eight faces, whose boundary is the boundary of the
missing face; so, since all edges are isomorphisms, the composite around the missing face is a
product of conjugates (by composites along edge paths) of the composites around the other eight
faces. Hence if eight of the nine faces commute, so does the ninth.

\emph{(4) Induction.} For $d\in I_v$ let $n(d)$ be the least $n$ for which a path as in the
statement exists. We prove $\mathrm P(a,b,c,d)$ for all labels $a,b,c$ by induction on $n(d)$. If
$n(d)=1$ then $d=v$, and this is the hypothesis together with (1). If $n(d')=n+1$, then
$N_{dv}^{d'}\ne0$ for some $d$ with $n(d)=n$, so $d'$ occurs in $d\otimes v$. Apply (3) to
$(a,b,c,d,v)$. The face $\mathrm P(a,b,c,d)\otimes\id_v$ commutes by induction; the faces
$\id_a\otimes\mathrm P(b,c,d,v)$, $\mathrm P(a\otimes b,c,d,v)$, $\mathrm P(a,b\otimes c,d,v)$ and
$\mathrm P(a,b,c\otimes d,v)$ have fourth object $v$ and commute by the hypothesis and (2); the
three squares commute. Hence $\mathrm P(a,b,c,d\otimes v)$ holds, and by (2) so does
$\mathrm P(a,b,c,d')$. By (1) this is the claim.
\end{proof}

For $\CF$ we apply Lemma~\ref{lem:reduction} with $v=V$: $V$ is self-dual and every simple
object is a summand of $V^{\otimes n}$, $n\le4$, so $I_V$ is everything, and the pentagon equation
for all $1{,}030{,}051$ five-tuples reduces to the $63{,}841$ five-tuples $(A,B,C,V,E)$.

\section{The construction}\label{sec:construction}
The construction is a computation over a finite field $\F_p$ with $p\equiv1\pmod{52}$ and a chosen
primitive $52$nd root of unity $q\in\F_p$, repeated for several primes and roots and followed by
exact reconstruction (Section~\ref{sec:reconstruction}). Modulo each prime and root:
\begin{enumerate}[label=(\arabic*)]
\item the cells: the action of the braiding, cup, cap and trivalent vertex of $V$ on two-step
paths, from explicit $\Uq$-modules (Sections~\ref{sec:modules}--\ref{sec:cells});
\item the simple objects as idempotents in $\End(V^{\otimes n})$, and an auxiliary basis of every
multiplicity space (Sections~\ref{sec:idempotents}, \ref{sec:gaugefix});
\item the F-symbols with third leg $V$ in the auxiliary bases (Section~\ref{sec:H});
\item the normal form of these F-symbols, defined by an explicit pivot rule
(Section~\ref{sec:nfrule});
\item all other F-symbols, by the pentagon recursion, run once in the normal-form gauge
(Section~\ref{sec:recursion});
\item the R-symbols, in the same gauge (Section~\ref{sec:Rsym}).
\end{enumerate}
Modulo each prime the computation passes through the auxiliary bases of step (2); they serve only to
produce the F-symbols with third leg $V$ modulo $p$, and no exact data in these bases are
reconstructed or used. The programs that implement each step are listed in Table~\ref{tab:scripts}.

\subsection{Modules over finite fields}\label{sec:modules}
We use the quantum group $\Uq$ with generators $E_i,F_i,K_i^{\pm1}$ in the conventions of
Jantzen's book~\cite{Jantzen} (we state them explicitly): $K_iE_jK_i^{-1}=q^{(\alpha_i,\alpha_j)}E_j$,
$[E_i,F_j]=\delta_{ij}(K_i-K_i^{-1})/(q_i-q_i^{-1})$ with $q_i=q^{(\alpha_i,\alpha_i)/2}$,
\[
\Delta(E_i)=E_i\otimes1+K_i\otimes E_i,\qquad \Delta(F_i)=F_i\otimes K_i^{-1}+1\otimes F_i,
\]
over $\F_p$, where $q$ is a primitive $52$nd root of unity in $\F_p$. Divided powers are not needed:
in all modules we construct, every $\alpha_i$-string of weights has length less than the order of
$q_i^2$ ($13$ for long and $26$ for short roots), so the quantum integers that occur are invertible.

A module $L(\lambda)$ is built from the top, weight space by weight space. If bases
of the weight spaces $L_{\nu+\alpha_j}$ are known, together with the matrices of all $E_i$ between
them, the candidates for $L_\nu$ are the vectors $F_jb$ with $b$ a basis vector of
$L_{\nu+\alpha_j}$. The matrix of each $E_i$ on a candidate $F_jb$ follows from the commutation
relation, and a linear combination of candidates is zero in the simple module if and only if it
is annihilated by every $E_i$. So $L_\nu$ is the quotient of the span of the candidates by the
kernel of the stacked matrix of the $E_i$, and we take as basis the pivot candidates of its row
echelon form. This constructs the simple quotient of the Verma module, truncated at a fixed depth
below the top weight. The dimension of every weight space constructed is compared with
Freudenthal's formula for the classical character (see~\cite{Humphreys}), so the modules
are the Weyl modules, which are simple at this root for the weights we use (the closed alcove).
Every basis vector is $F_j$ applied to a basis vector one level up, which defines, for any module
$M$ and vector $m$ of weight $\lambda$ annihilated by the $E_i$, the $U^-$-module map
$L(\lambda)\to M$ sending the top vector to $m$: it is computed by applying the same words in the
$F_j$ to $m$.

\subsection{The path model}\label{sec:paths}
Let $\cC=\CF$ and $V=L(\ob{0001})$. For simple $A,B$ with $N_{AV}^B\ne0$ we call a basis vector
$e\in\Hom(B,A\otimes V)$ an \emph{edge} $A\to B$; there are $N_{AV}^B\in\{1,2\}$ of them
(Table~\ref{tab:graph}). For a simple object $X$ and $n\ge0$ let $P_X(n)$ be the vector space with
basis the edge paths
\[
X=Z_0\xrightarrow{e_1}Z_1\xrightarrow{e_2}\cdots\xrightarrow{e_n}Z_n ,
\]
and $P_X(T,n)$ the span of the paths ending at $T$. The path corresponds to the morphism
\[
(e_1\otimes\id_{V^{\otimes(n-1)}})\circ(e_2\otimes\id_{V^{\otimes(n-2)}})\circ\cdots\circ e_n
\in\Hom(T,X\otimes V^{\otimes n}),
\]
and by semisimplicity these morphisms form a basis of $\Hom(T,X\otimes V^{\otimes n})$. A
morphism $f:V^{\otimes n}\to V^{\otimes m}$ therefore acts on path spaces by
$\id_X\otimes f:P_X(T,n)\to P_X(T,m)$, for every base $X$ and end point $T$. The following
generators suffice: the braiding $\sigma_i$ of the strands $i,i+1$, the cup and cap of $V$
inserted or applied at position $i$, and the trivalent split $t:V\to V\otimes V$ and merge
$m:V\otimes V\to V$. Each acts \emph{locally}: by the interchange law, applying
$\id\otimes\sigma_i\otimes\id$ to a path changes only the two steps $Z_{i-1}\to Z_i\to Z_{i+1}$,
replacing the two-step path $Z_{i-1}\to Z_i\to Z_{i+1}$ by its image in
$\Hom(Z_{i+1},Z_{i-1}\otimes V\otimes V)$, which is a combination of two-step paths from
$Z_{i-1}$ to $Z_{i+1}$. Similarly the cup at a vertex $Z_i$ inserts a combination of two-step
paths $Z_i\to B\to Z_i$, the cap contracts a two-step path $Z_i\to B\to Z_i$ to a multiple of the
empty path at $Z_i$, the split replaces one step by a combination of two-step paths with the same
end points, and the merge does the opposite. So all generators act by sparse matrices determined
by the following finite data, which we call the \emph{cells}: for all simple $A,C$, with
$P(A\to C,2)$ the span of the two-step paths $A\to B\to C$,
\begin{itemize}
\item the matrix $\sigma^{(A,C)}$ of $\id_A\otimes c_{V,V}$ on $P(A\to C,2)$;
\item the vector $\mathrm{cup}^{(A)}\in P(A\to A,2)$ and the functional $\mathrm{cap}^{(A)}$ on
$P(A\to A,2)$;
\item the matrices of $\id_A\otimes t$ from edges $A\to C$ to $P(A\to C,2)$, and of $\id_A\otimes
  m$ in the other direction.
\end{itemize}
For $\CF$ these are $4{,}472$ numbers.

\subsection{Cells from the quantum group}\label{sec:cells}
To compute the cells we realise edges as highest weight vectors. An edge $e:L(B)\to L(A)\otimes
V$ is determined by the image of the top vector of $L(B)$, a vector of weight $B$ in
$L(A)\otimes V$ annihilated by all $E_i$. Two-step paths $A\to B\to C$ are vectors of weight $C$
in $L(A)\otimes V\otimes V$ annihilated by all $E_i$. Such vectors are computed from their
components along the top vector $v_A$ of $L(A)$: if $y$ is a
highest weight vector of weight $C$ in $L(A)\otimes W$, write $\tau(y)\in W_{C-A}$ for the
coefficient of $v_A$. Then $\tau$ is injective on highest weight vectors (the lower components are
forced by the equations $\Delta(E_i)y=0$, solved weight by weight downwards), and $\tau$ commutes
with every operator acting on $W$ alone. For $W=V\otimes V$ this reduces all computations to the
$676$-dimensional space $V\otimes V$, in fact to its weight space $(V\otimes V)_{C-A}$.

The image $\tau\bigl((e_1\otimes\id_V)\circ e_2\bigr)$ of a two-step path is computed as follows.
Since $\Delta(F_j)=F_j\otimes K_j^{-1}+1\otimes F_j$ and $F_j$ lowers weights in the first factor,
$\tau(F_jx)=F_j\tau(x)$ for $x\in L(A)\otimes V$; hence $\tau\circ e_1:L(B)\to V$ is the
$U^-$-module map sending the top vector of $L(B)$ to $\tau(e_1)$, computed by the words in the
$F_j$ that define the basis of $L(B)$ (Section~\ref{sec:modules}). Then
$\tau((e_1\otimes\id)e_2)=((\tau\circ e_1)\otimes\id_V)(e_2)$, with $e_2$ the highest weight
vector of weight $C$ in $L(B)\otimes V$.

\emph{Negligible morphisms.} The constituents of $L(A)\otimes V$ lie in the closed alcove, since
every weight of $V$ has level at most $1$; those of level $5$ lie on the wall $(\lambda+\rho,\theta)=26$
and have quantum dimension zero. By the linkage principle $L(A)\otimes V$ is the direct sum of its
constituents. Hence the space of highest weight vectors of weight $C$ in $L(A)\otimes V\otimes V$
is spanned by the two-step paths through $B$ in the open alcove together with the vectors
$(e\otimes\id_V)\circ f$, where $e:L(B)\to L(A)\otimes V$ with $B$ on the wall and $f$ runs over
the highest weight vectors of weight $C$ in $L(B)\otimes V$. The second kind of vector is a
morphism $L(C)\to L(A)\otimes V\otimes V$ that factors through the negligible object $L(B)$, so it
is negligible. We compute both kinds of vectors through $\tau$ and check that together they have
rank equal to the classical multiplicity of $L(C)$ in $L(A)\otimes V\otimes V$, for every pair
$(A,C)$. To apply an operator $\id_A\otimes f$ ($f$ an endomorphism of $V\otimes V$, or a map from
$V$ or to $V$) to a path, we apply $f$ to its $\tau$-image in $V\otimes V$, expand the result in
this basis, and discard the coefficients of the negligible vectors. This is the action of $f$ in
the semisimple quotient, in the path basis.

\emph{The generators on $V\otimes V$.} The module $V\otimes V$ is semisimple at $q$ (its five
constituents lie in the open alcove), and we compute its decomposition explicitly. The braiding
is
\[
c_{V,V}=\sum_W\epsilon_W\,q^{((W,W+2\rho)-2(V,V+2\rho))/2}\,\pi_W
\]
with $\pi_W$ the projection on the summand $W$ and $\epsilon_W=+1$ on the summands of $S^2V$
($\one$, $V$, $L(\ob{0002})$) and $-1$ on those of $\Lambda^2V$ ($L(\ob{0010})$, $L(\ob{1000})$);
the eigenvalues are $q^{-24},q^{-12},q^2,-1,-q^{-6}$ respectively. This is the spectral form of
the braiding of $\Uq$ on $V\otimes V$ (its square on the summand $W$ is $\theta_W/\theta_V^2$);
we check the Yang--Baxter equation on $V^{\otimes3}$. The cup is the invariant vector of
$V\otimes V$ and the cap the invariant functional with $\mathrm{cap}(\mathrm{cup})=1$; the
zigzag composite is then a scalar $z$, and we rescale the cap by $z^{-1}$, so that the zigzag is
the identity and the loop is $z^{-1}$. The split $t$ is the inclusion of the summand $V$ of
$V\otimes V$ (determined by its highest weight vector) and the merge $m$ is the projection with
$mt=\id_V$.

\emph{Gauge of the edges.} The basis of $\Hom(B,A\otimes V)$ is chosen so that the
$\tau$-images, which lie in $V_{B-A}$ (of dimension at most $2$), form a reduced echelon basis.

\emph{Checks.} At every base $A$ we verify the braid relation on
$P_A(3)$, the minimal polynomial of $\sigma$, both zigzag identities, naturality of the merge
and of the cup with respect to the braiding, and the theta graph. The curl is $q^{24}=\theta_V$
and the loop is the quantum dimension $d_V$. The whole step takes about $10$ seconds.

\subsection{Simple objects as idempotents}\label{sec:idempotents}
For each simple $S$ of depth $s$ we choose a \emph{spectral path} $\gamma_S=(\one=Z_0,V=Z_1,
Z_2,\dots,Z_s=S)$ in the fusion graph of $V$ such that every step $Z_{k-1}\to Z_k$ has
multiplicity one and $\theta_{Z_k}$ differs from the twists of all other constituents of
$Z_{k-1}\otimes V$. Let
$\Theta_k=\theta_V^k(\sigma_1\sigma_2\cdots\sigma_{k-1})^k$ be the ribbon twist of $V^{\otimes k}$
(the full twist braid times $\theta_V^{\otimes k}$), acting on the first $k$ strands; it acts on a
summand $Z$ of $V^{\otimes k}$ by $\theta_Z$. Define
\[
p_S=\prod_{k=2}^{s}\ \prod_{t}\frac{\Theta_k-t}{\theta_{Z_k}-t},
\]
the inner product over the distinct twists $t\ne\theta_{Z_k}$ of the other constituents of
$Z_{k-1}\otimes V$, the factors for smaller $k$ applied first. This is an idempotent in
$\End(V^{\otimes s})$ with image a copy of $S$: on the image of $p_{Z_{k-1}}\otimes\id_V$, the
operator $\Theta_k$ acts on the summand $B$ of $Z_{k-1}\otimes V$ by $\theta_B$, and the Lagrange
factor keeps the summand $Z_k$. At the base $\one$, $p_S$ is the projection on the single path
$\gamma_S$. One exception is needed: $\theta_{\ob{0002}}=\theta_\one=1$, so at the second step
we first remove the unit summand of $V\otimes V$ with $1-\mathrm{cup}\circ\mathrm{cap}$ (with the
raw cap, $\mathrm{cap}(\mathrm{cup})=1$). With this, every spectral path has length equal to the
depth. We write $\iota_S:S\to V^{\otimes s}$ and $\pi_S:V^{\otimes s}\to S$ for the inclusion and
projection with $\iota_S\pi_S=p_S$.

\subsection{Auxiliary bases of the multiplicity spaces}\label{sec:gaugefix}
For simple $X,S,T$, the space $\Hom(T,X\otimes S)$ is identified with the image of
$\id_X\otimes p_S$ on $P_X(T,s)$, a subspace of dimension $N_{XS}^T$ of the space spanned by the
paths of length $s$ from $X$ to $T$. As auxiliary basis we take its reduced echelon basis with respect
to the following order of paths: by end point (in the order of Table~\ref{tab:objects}), then
lexicographically by the sequence of vertices (as Dynkin label tuples) and then of edge indices.
These are the splitting maps $\psi^{XS}_{T,l}$, $0\le l<N_{XS}^T$, used in steps (3) and (6). The
basis depends only on the subspace and the order; to find it we apply $\id_X\otimes p_S$ to
$N_{XS}^T+3$ random vectors and check that the rank is $N_{XS}^T$. All $2{,}302$ images have the
fusion dimension, so the construction reproduces the fusion rules.

\subsection{F-symbols with third leg \texorpdfstring{$V$}{V}}\label{sec:H}
In the \emph{splitting convention} the associator is described by matrices $G$ with
\[
(\psi^{XY}_{R,i}\otimes\id_W)\circ\psi^{RW}_{T,j}=\sum_{S,k,l}G[X,Y,W,T]_{(R,i,j),(S,k,l)}\,
(\id_X\otimes\psi^{YW}_{S,k})\circ\psi^{XS}_{T,l}
\]
(both sides in $\Hom(T,X\otimes Y\otimes W)$). For $W=V$ both bases can be written in the path
model at base $X$. The left basis vectors are concatenations: the path
$\psi^{XY}_{R,i}\in P_X(R,y)$ followed by the edge $R\to T$ given by $\psi^{RV}_{T,j}$. For the
right basis we need the morphism $\id_X\otimes\psi^{YV}_{S,k}$, realised on tensor powers of $V$
as
\[
\Phi_k=(\iota_Y\otimes\id_V)\circ\psi^{YV}_{S,k}\circ\pi_S:V^{\otimes s}\to V^{\otimes(y+1)} .
\]
It is characterised by its value at the base $\one$, where $\pi_S$ maps the path $\gamma_S$ to
the generator of $S$ and $\Phi_k\gamma_S$ is the path $\gamma_Y$ followed by the $k$-th edge
$Y\to S$. To realise $\Phi_k$ at an arbitrary base we take random words $w$ in the generators
(braidings, splits, cups and merges) from $V^{\otimes s}$ to $V^{\otimes(y+1)}$, compress them to
$(p_Y\otimes\id_V)\circ w\circ p_S$, which lie in the $N_{YV}^S$-dimensional space
$\Hom(S,Y\otimes V)$, evaluate them at the base $\one$ until $N_{YV}^S$ independent ones are
found, and take the linear combination that maps $\gamma_S$ to $\gamma_Y$ followed by the $k$-th
edge. Applying this combination at the base $X$ to the basis $\psi^{XS}_{T,l}$ gives the right
basis vectors, and $G[X,Y,V,T]$ is obtained by solving a linear system. The words only serve to
span $\Hom(S,Y\otimes V)$, and the result does not depend on them (two different random seeds give
identical results). There are $3{,}451$ such matrices, with $399{,}516$ entries, the largest of size $78$. We write
$F[X,Y,V,T]=(G[X,Y,V,T]^{-1})^{\mathsf T}$ for the same matrices in the fusion-map convention of
Section~\ref{sec:conventions} (see Section~\ref{sec:convchange}).

\subsection{The normal form}\label{sec:nfrule}
A gauge transformation (Section~\ref{sec:conventions}) is a family of invertible matrices
$g^{XY}_Z$; we only use unit-preserving ones, $g^{XY}_Z=1$ when $X$ or $Y$ is the unit. The
\emph{normal form} is the gauge transformation $g$ determined from the F-symbols with third leg $V$ by
the following rule, and $\mathrm{NF}(F)=g\cdot F$. All label sets are ordered by (position of the
object in Table~\ref{tab:objects}, then the multiplicity indices).
\begin{enumerate}[label=(\arabic*),leftmargin=2em]
\item \emph{Depth and parents.} The depth of $S$ is its distance from $\one$ in the fusion graph of
$-\otimes V$. For depth at least $2$, the parent $P(S)$ is the first object $P$ of depth one less with
$N_{PV}^S=1$.
\item \emph{One-dimensional edges.} The unknowns are the scalars $x_t=g^{AV}_C$ for the spaces
$t=(A,V,C)$ with $A\ne\one$ and $N_{AV}^C=1$. Run over the blocks $F[A,V,V,C]$, $A\ne\one$, in order
of $(A,C)$, over the columns $(Z',0,0)$ with $Z'\in\{\one,V\}$ (and $N_{AV}^C=1$ if $Z'=V$), and over
the rows $(Z,0,0)$ with $N_{AV}^Z=N_{ZV}^C=1$, keeping the nonzero entries. By~\eqref{eq:gauge-scalar}
each such entry is multiplied by a monomial in the unknowns, with integer exponent vector $D_e$.
Reduce each $D_e$ in turn by the pivot rows found so far; if the result is nonzero it must contain a
coefficient $\pm1$, and its first such variable becomes a new pivot. The gauge is $x_t=1$ for the
non-pivot variables and, for the pivot variables, the values (solved in reverse order) that make every
pivot entry equal to $1$.
\item \emph{Two-dimensional loops} $\Hom(A\otimes V,A)$, for $A=\ob{0011},\ob{0012}$. For one-dimensional
neighbours $C$ the vectors $(F[A,V,V,C]_{(A,i,0),(V,0,0)})_{i=0,1}$, divided by $g^{VV}_V$, transform by
the $2\times2$ gauge matrix of the loop; for the first pair of such vectors (in the order of $C$)
with nonzero determinant, the gauge matrix is the inverse of the matrix with these two columns, so
that the two vectors become the standard basis vectors.
\item \emph{Second legs of depth at least $2$}, in order of depth, then of $S$, then of $A\ne\one$,
then of $C$: let $M$ be the block of columns $(S,0,l)$ of $F[A,P(S),V,C]$, with the gauges fixed so far
applied and divided by $g^{P(S)V}_S$. Its pivot rows are the lexicographically first maximal set of
linearly independent rows, and $g^{AS}_C$ is the submatrix of $M$ on these rows; since $g^{AS}_C$
acts on these columns by its inverse (Section~\ref{sec:conventions}), the normalised block is the
identity on the pivot rows.
\end{enumerate}
Every choice in this rule depends only on the gauge class: on which entries of blocks with
one-dimensional multiplicity spaces vanish, on linear independence of vectors that are already
normalised, and on exponents. So $\mathrm{NF}(h\cdot F)=\mathrm{NF}(F)$ for every gauge $h$, and,
since the rule uses only field operations, $\mathrm{NF}(\sigma\cdot F)=\sigma\cdot\mathrm{NF}(F)$ for
every field automorphism $\sigma$ (Section~\ref{sec:fields}). The two requirements of the rule (a unit
coefficient in every nonzero reduced exponent vector; enough independent vectors and rows) are met
for $\CF$. In step (2) there are $89$ unknowns and $467$ entries, the exponent vectors have rank $74$,
and all $74$ pivots are units, so no roots are needed. The index choices are computed at a reference
prime and stored; at every prime and root the rule is recomputed and asserted to give the same
choices (a prime at which a relevant minor vanishes would give different choices and is not used).
The choices therefore satisfy the nondegeneracy conditions over the field, and the rule with these
fixed choices is a function of the gauge class; this is all that is used below. At every prime tested they are also the
greedy choices.

\subsection{The pentagon recursion}\label{sec:recursion}
All other F-symbols follow from those with third leg $V$ by the pentagon
equation~\eqref{eq:pentagon-matrix}. The recursion is applied once, to the normal form of the
F-symbols with third leg $V$ (in the splitting convention, $G=(\mathrm{NF}(F)^{-1})^{\mathsf T}$),
so that its output is the whole table in the normal-form gauge. The pentagon equation has the same
shape $M_1M_2=N_1N_2N_3$ for the $G$-blocks. For each simple $S\ne\one$
let $S'$ be the predecessor of $S$ on its spectral path, so that $S$ is a summand of $S'\otimes V$
with multiplicity one.
\begin{enumerate}[label=Step \arabic*.,leftmargin=4.5em]
\item (Second leg $V$.) For the five-tuple $(X,V,W',V,E)$, solve
$M_2=M_1^{-1}N_1N_2N_3$. The blocks of $M_2$ are $G[X,V,R,E]$ for all $R$ in $W'\otimes V$, and
the factors on the right involve only F-symbols with third leg $V$ and $G[X,V,W',\cdot]$, of
smaller depth of the third leg. By induction on the depth of $W$ this gives all $G[X,V,W,E]$.
\item (General.) For the five-tuple $(X,Y',V,W,E)$, solve $N_2=N_1^{-1}M_1M_2N_3^{-1}$. The blocks
of $N_2$ are $G[X,T,W,E]$ for all $T$ in $Y'\otimes V$, and the other factors involve F-symbols
with third leg $V$, second leg $V$ (Step 1), and second leg $Y'$ of smaller depth. By induction on
the depth of the second leg this gives everything.
\end{enumerate}
Only the rows of the products that are needed are formed. The F-matrices with a unit leg are set
to identities. Each step produces the blocks for every $R$ (respectively $T$) in $W'\otimes V$
(respectively $Y'\otimes V$) and every multiplicity copy, so it checks itself: the off-diagonal
blocks must vanish, the copies must agree, blocks already known (from smaller depth, the unit legs
or the family with third leg $V$) must be reproduced, and blocks of the same depth reached from
different predecessors must agree; there were no failures. For $G[X,\cdot,
\cdot,\cdot]$ Step 2 uses only blocks with the same first leg $X$ besides the families with second or
third leg $V$, so the table can be computed one first leg at a time; with this, all $58{,}740$
matrices take about $640$ seconds on one core per prime and root of unity, after about $150$
seconds for the cells, the family with third leg $V$ and its normal form.

The recursion uses only pentagons, so its output is determined by its input:

\begin{lemma}\label{lem:Hdetermines}
Let $F_1,F_2$ be two families of invertible F-matrices for the same fusion rules, both
satisfying the pentagon equation, with identity F-matrices on unit legs, such that every simple
object is a summand of a tensor power of $V$. If $F_1$ and $F_2$ agree on all blocks
$[X,Y,V,T]$, they agree everywhere.
\end{lemma}

\begin{proof}
By induction on the depth of the third leg $W$, the cases of depth $0$ and $1$ being the
hypothesis. For $W$ a summand of $W'\otimes V$, the pentagon equation for $(X,Y,W',V,E)$ reads
$M_1M_2=N_1N_2N_3$, where $M_1$ consists of blocks $[P,W',V,E]$ and $N_1,N_2,N_3$ of blocks with
third leg $W'$ or $V$. So $M_2=M_1^{-1}N_1N_2N_3$ is determined, and its blocks are the
$F[X,Y,R,E]$ for all $R$ in $W'\otimes V$.
\end{proof}

\subsection{R-symbols}\label{sec:Rsym}
In the splitting convention, $c_{X,Y}\circ\psi^{XY}_{Z,i}=\sum_jR^{\mathrm s}[X,Y,Z]_{i,j}\,
\psi^{YX}_{Z,j}$. At the base $\one$, the vector $(\iota_X\otimes\iota_Y)\circ\psi^{XY}_{Z,i}$ is
the path $\gamma_X$ followed by the $i$-th basis vector of $P_X(Z,y)$, and $c_{X,Y}$ corresponds
to the positive braid $c_{V^{\otimes x},V^{\otimes y}}$ on $x+y$ strands. So $R^{\mathrm s}$ is
obtained by applying a braid word to $N_{XY}^Z$ path vectors and expressing the results in the
paths $\gamma_Y$ followed by the basis of $P_Y(Z,x)$. This gives the R-symbols in the auxiliary bases;
the normal-form gauge $g$ of Section~\ref{sec:nfrule} is then applied modulo $p$,
$R[X,Y,Z]\mapsto g^{YX}_ZR[X,Y,Z](g^{XY}_Z)^{-1}$ in the fusion-map convention. Since $g$ is real, the
same is done at the conjugate root $q^{-1}$ (cells and R only).

\subsection{Change of convention}\label{sec:convchange}
The fusion maps of Section~\ref{sec:conventions} are the dual bases of the splitting maps,
$\phi^{XY\to Z}_i\circ\psi^{XY}_{Z,j}=\delta_{ij}\id_Z$. Then
\[
F[X,Y,W,T]=\bigl(G[X,Y,W,T]^{-1}\bigr)^{\mathsf T},\qquad R[X,Y,Z]=R^{\mathrm s}[X,Y,Z]^{\mathsf T}.
\]
The normal form and all data files and statements in this paper refer to the fusion-map
convention.

\subsection{Dimensions and twists}
The trace of $p_X$ at the base $\one$, computed with nested cups and caps, is the quantum
dimension $d_X$ for every $X$, and the twists recomputed from the R-symbols are
$q^{(X,X+2\rho)}$, both modulo primes. So the pivotal structure realised by the cup and cap
of $V$ is the spherical structure with positive dimensions at the unitary root.

\section{Fields}\label{sec:fields}\label{sec:normalform}
\subsection{Galois action}
Every step of Section~\ref{sec:construction} is a field operation on quantities defined over
$\Q(q)=\K$, with choices fixed by the conventions (reduced echelon forms, fixed orders, the stored
choices of the normal form). So for every embedding $\K\to\F_p$, $\zeta\mapsto z$, the computation at
$q=z$ gives the reduction of the same elements of $\K$, provided no pivot or denominator vanishes
modulo the prime, which is checked. The $24$ primitive $52$nd roots of unity $z^w$,
$w\in(\Z/52)^\times$, give the $24$ embeddings, and the output at $q=z^w$ is the image of the exact
data under $\sigma_w:\zeta\mapsto\zeta^w$, the data of the Galois conjugate category.

\subsection{Normal forms and the field of moduli}
Write $\sigma\cdot F$ for the Galois conjugate of a family of F-symbols and $h\cdot F$ for its
gauge transform. A \emph{normal form} is a rule $F\mapsto\mathrm{NF}(F)$ such that
$\mathrm{NF}(F)$ is gauge equivalent to $F$, $\mathrm{NF}(h\cdot F)=\mathrm{NF}(F)$ for every gauge
$h$, and $\mathrm{NF}(\sigma\cdot F)=\sigma\cdot\mathrm{NF}(F)$ for every field automorphism
$\sigma$; the last property holds for any rule built from field operations with index choices
fixed in advance. If $\sigma\cdot F$ is gauge equivalent to $F$, then
$\sigma\cdot\mathrm{NF}(F)=\mathrm{NF}(\sigma\cdot F)=\mathrm{NF}(F)$. So the field generated by
the entries of $\mathrm{NF}(F)$ is contained in the fixed field of every such $\sigma$, that is in
the \emph{field of moduli}, and conversely $\mathrm{NF}(F)$ is a complete invariant of the gauge
class. A normal form therefore gives the field of moduli and a gauge realising it at the same time;
in particular no cohomological obstruction to descent arises.

\subsection{The fields of the normal form of \texorpdfstring{$\CF$}{CF}}
At every prime used, the residues of the normal form at the roots $z^w$ and $z^{25w}$ agree, for
every entry of the table (the residues of the F-symbols in the auxiliary bases differ at these roots,
for example in $192{,}713$ of the entries with third leg $V$ at one prime). So the normal form is fixed by
$\sigma_{25}$. It is also fixed by complex conjugation $\sigma_{51}$: the F-symbols are real (their
residues at $z^w$ and $z^{-w}$ agree; checked on all entries at one prime). Hence it is fixed by
$\sigma_{27}$, and its F-symbols lie
in $\kk=\Kt^+$, the fixed field of $\{1,25,27,51\}$, of degree $6$, and are reconstructed directly over
$\kk$ (Section~\ref{sec:reconstruction}). The R-symbols in the same gauge are fixed by $\sigma_{27}$
and lie in $\Kt=\Q(\zeta^4)$, of degree $12$. The minimal polynomial of $\eta_1=\zeta_{13}+\zeta_{13}^{-1}$
is $x^6+x^5-5x^4-4x^3+6x^2+3x-1$.

\emph{Minimality, F.} On the reconstructed F-symbols with third leg $V$, the automorphisms $\sigma_a$
that fix every entry are exactly $a\in\{1,25,27,51\}$ (\texttt{nf\_galois.py}). Since $\mathrm{NF}$ is a
complete invariant of the gauge class, if some gauge $h\cdot F$ had all entries in a subfield
$E\subseteq\kk$, every $\sigma$ fixing $E$ would satisfy $\sigma\cdot F\sim F$, hence would fix
$\mathrm{NF}(F)$; so $E=\kk$. With the labels of the objects fixed, $\kk$ is the field of moduli and
the field of definition of the fusion category.

\emph{Minimality, F and R.} The one-dimensional $R[X,X,Z]$ are gauge invariant, and the only
$a\in(\Z/52)^\times$ with $\sigma_a$ fixing all $100$ of them are $a=1,27$ (\texttt{autoeq\_checks.py});
so no gauge has the braided data in a proper subfield of $\Q(\zeta)^{\langle\sigma_{27}\rangle}=\Kt$.
Every one-dimensional $R[X,X,Z]$ is $q^e$ with $e$ even, and the characteristic polynomials of all
$88$ matrices $R[X,X,Z]$ of size at least $2$ have coefficients in $\Kt$ (\texttt{r\_invariants.py}).

\emph{Comparison.} Galindo~\cite[Theorem~3.7]{Galindo} proves that the fusion
category underlying $\cC(\mathfrak g,k)$ has a split form over $\Q(\zeta+\zeta^{-1})$ with
$\zeta=e^{\pi i/(r^\vee(k+h^\vee))}$, $r^\vee$ the lacing number; for $\CF$ this is $\Lf=\Q(\zeta_{52})^+$,
of degree $12$. The normal form gives a split form over the subfield
$\kk$ of index $2$, and by the argument above no smaller field is possible with the labels fixed. For
the braided category, Galindo shows that no split real form exists for the standard braiding except
for $\cC(E_8,1)$ and $\cC(D_{4m},1)$; this agrees with $\Kt$, which is not real, being the smallest
field for the braided data.

\emph{$q$ and $-q$.} Since $\sigma_{27}$ ($q\mapsto-q$) fixes $\kk$ and $\Kt$, the normal-form data at
$q$ and at $-q$ are identical: the braided categories at $q$ and $-q$ are equivalent, with the labels of
the objects fixed.

\section{Exact reconstruction}\label{sec:reconstruction}
\subsection{Primes and roots}
We use primes $p\equiv1\pmod{52}$ just below $2^{30}$. For each prime, $z=g^{(p-1)/52}$ with $g$ the
least primitive root modulo $p$, and the run at $(p,w)$ uses $q=z^w$. An element of $\kk$ is
determined by its values at one root from each of its $6$ embeddings, $w=1,3,5,7,9,11$ (the root
$z^w$ and its images under $\{1,25,27,51\}$ give the same value), and an element of $\Kt$ by its values
at the $12$ roots $w=1,3,\dots,11$ and $41,43,\dots,51$, which are distinct modulo $13$. The residues of
the coefficients in the bases $\eta_1,\dots,\eta_6$ of $\kk$ ($\eta_j=\zeta_{13}^j+\zeta_{13}^{-j}$) and
$1,\zeta_{13},\dots,\zeta_{13}^{11}$ of $\Kt$ follow by solving Vandermonde-type systems.

\subsection{Reconstruction}
For each coefficient, the residues at several primes are combined by the Chinese remainder theorem
and a rational number $a/b$ with $|a|,b\le\sqrt{M/2}$ is recovered by rational
reconstruction~\cite{WangGuyDavenport,MCA}, $M$ the product of the primes. Further primes are held out
and the reconstructed values are compared with the residues computed there.
\begin{itemize}
\item \emph{F-symbols}: $9$ primes at the $6$ roots, and $2$ held-out primes at all $12$ real roots
($78$ runs). Every coefficient reconstructs; every value agrees with the held-out primes, whose
residues at $z^w$ and $z^{25w}$ also agree in every entry. The table has $328{,}727{,}941$ nonzero
coefficients in the basis $\eta_1,\dots,\eta_6$, with numerators and denominators of at most $99$ bits.
A preliminary probe of $200{,}000$ entries at $20$ primes found that $7$ primes suffice for all of them.
\item \emph{R-symbols and their inverses}: over $\Kt$, at the $12$ roots above, with heights at most
$47$ and $41$ bits.
\end{itemize}
The exact tables in the archive were reconstructed in this way from residues of the normal form
computed modulo each prime.

The residues from the pipeline of Section~\ref{sec:construction}, run at all $78$ (prime, root) pairs,
reconstruct to files identical byte for byte to the archived ones (\texttt{deposit/logs/finish\_all.log}).

Rational reconstruction from finitely many primes is a heuristic step: it cannot by itself certify
that the fractions found are the values defined by the construction. Theorem~\ref{thm:main} is about
the reconstructed data, and its proof in Section~\ref{sec:verification} uses only the data files.

\section{Verification}\label{sec:verification}
All checks in this section are run on the data files themselves: the F-table over $\kk$ (stored in
the real basis of $\Lf$, so that the same verifiers read it) and the R-symbols over $\Kt$. The
verifying programs share no code with the programs that produced the data, apart from the reader of
the binary format.

\subsection{Height bounds}
Write $\Z[\zeta]=\Z[x]/(\Phi_{52})$, with
$\Phi_{52}(x)=\Phi_{13}(-x^2)=\sum_{k=0}^{12}(-1)^kx^{2k}$, and for $\rho\in\Z[\zeta]$ let
$\norm\rho$ be the largest absolute value of a coefficient of its representative of degree less
than $24$. For $0\le k\le46$ write $\zeta^k=\sum_{j<24}r_{kj}\zeta^j$ and put
$c=\max_j\sum_{k=0}^{46}(\min(k,46-k)+1)\,|r_{kj}|$; a direct computation gives $c=46$. Then
$\norm{\rho\sigma}\le c\norm\rho\norm\sigma$: the product of the representatives has degree at most
$46$, its coefficient of $x^k$ is a sum of at most $\min(k,46-k)+1$ products, and reducing $x^k$ to
$\sum_jr_{kj}x^j$ gives the bound. An element
$a_0+\sum_{k=1}^{11}a_kc_k$ with $c_k=\zeta^k+\zeta^{-k}$ and integer $a_k$ has $\norm\cdot\le2\max|a_k|$ in the power basis.

\begin{lemma}\label{lem:height}
Let $\rho\in\Z[\zeta]$ and let $\mathcal P$ be a set of primes $p\equiv1\pmod{52}$ such that the
image of $\rho$ under $\zeta\mapsto z$ vanishes for each of the $24$ roots $z$ of $\Phi_{52}$ in
$\F_p$, for every $p\in\mathcal P$. If $\norm\rho<\prod_{p\in\mathcal P}p$ then $\rho=0$. If $\rho$
is real (fixed by $\zeta\mapsto\zeta^{-1}$), it suffices that the images vanish at $12$ roots, one
from each pair $\{z,z^{-1}\}$.
\end{lemma}

\begin{proof}
The reduction of the representative of $\rho$ modulo $p$ has degree less than $24$ and $24$
distinct roots, so all its coefficients are divisible by $p$, hence by $\prod p$, hence zero. If
$\rho$ is real, its image at $z^{-1}$ is the image of $\bar\rho=\rho$ at $z$.
\end{proof}

\begin{lemma}\label{lem:matrix-bound}
Let $A_1,\dots,A_r$ and $B_1,\dots,B_t$ be matrices over $\K$ such that the products
$A=A_1\cdots A_r$ and $B=B_1\cdots B_t$ are defined and of the same shape, with inner dimensions
at most $n$. Let $s_M$ be positive integers with $s_MM$ integral over $\Z[\zeta]$ and
$\norm{s_MM_{ij}}\le H_M$ for all entries. If for every $p$ in a set $\mathcal P$ of primes, none
dividing any $s_M$, the reductions satisfy $A=B$ at all $24$ roots, and
\[
\prod_{p\in\mathcal P}p>n^{r-1}c^{r-1}\prod_iH_{A_i}\prod_js_{B_j}+n^{t-1}c^{t-1}\prod_jH_{B_j}\prod_is_{A_i},
\]
then $A=B$ over $\K$.
\end{lemma}

\begin{proof}
Multiply $A-B$ by $\prod s_{A_i}\prod s_{B_j}$; each entry becomes an element of $\Z[\zeta]$
whose norm is bounded by the right hand side, since an entry of a product of $r$ matrices is a sum of
at most $n^{r-1}$ products of $r$ entries. Apply Lemma~\ref{lem:height}.
\end{proof}

In the programs, $s_M$ is the least common multiple of the common denominators of the blocks
assembled into $M$ (in the pentagon, for the scalar equation in a given row and column, only the
blocks of $M_1,M_2,N_1,N_3$ that meet that row or column), $H_M$ is computed from the stored
numerators with the factor $2$ for the change of basis, and the primes are the primes
$p\equiv1\pmod{52}$ below $2^{31}$, in decreasing order, that divide no denominator; each identity
is checked modulo the first $k$ of them, with $k$ least such that the product exceeds twice the
bound. The checks are done with full matrices, independently for each four-tuple or five-tuple.
The verifying programs read only the exported data files and share no code with the programs
that produced the data, apart from the reader of the binary format.

\subsection{Hexagons}
\texttt{verify\_hex\_bin.py} checks both identities~\eqref{eq:hex-matrix} for all $58{,}740$
four-tuples at all $24$ roots, and $RR^{-1}=1$ for all R-blocks. The bounds (twice the bound of
Lemma~\ref{lem:matrix-bound}) reach $1293$ bits, and the number of primes needed per four-tuple
ranges from $1$ to $42$. There were no failures. Since every F-matrix is invertible (below), this
proves Theorem~\ref{thm:main}(i). The computation took $8593$ seconds with $2$ workers. We also
checked $RR^{-1}=1$ directly in exact rational arithmetic (\texttt{exact\_checks.py}).

\subsection{Invertibility, units and the dimensions}\label{sec:invert}
All $58{,}740$ F-matrices are invertible modulo $p=1073741189$, hence over $\K$ (the determinant of
the scaled matrix is an element of $\Z[\zeta]$ with nonzero image). The R-matrices are invertible
since $RR^{-1}=1$. The following are checked exactly on the data (\texttt{exact\_checks.py}): the
$6{,}859$ F-matrices with a unit among their first three legs are identities in the sense of
Section~\ref{sec:conventions}; $N_{XX}^\one=1$ for all $X$; and
$F[X,X,X,X]_{(\one,0,0),(\one,0,0)}\cdot d_X=1$ in $\K$ for all $16$ objects, with $d_X$ given
by~\eqref{eq:qdim}. This proves Theorem~\ref{thm:main}(iii).

\subsection{The pentagon equation}\label{sec:pentagon}
The pentagon equation~\eqref{eq:pentagon-matrix} for all five-tuples amounts to about
$4\times10^{11}$ scalar equations, with matrices up to $19{,}900\times19{,}900$. It follows from the
pentagons with fourth object $V$ and Lemma~\ref{lem:reduction}.

\emph{Fourth object $V$} (\texttt{pent\_reduced.py}). For all $63{,}841$ five-tuples $(A,B,C,V,E)$ (the
largest of dimension $n_{ABCVE}=3380$; $\sum n^3\approx1.8\times10^{12}$), the
identity~\eqref{eq:pentagon-matrix} is checked as full matrices at the $12$ real roots
(Lemma~\ref{lem:height}; the F-symbols are real) of the first $k$ primes $p\equiv1\pmod{52}$ below
$2^{31}$, with $k$ the least number for which the product exceeds twice the bound of
Lemma~\ref{lem:matrix-bound} for that five-tuple. The bounds reach $1176$ bits; between $1$ and $38$
primes are needed (mean $9.1$). There were no failures, so the pentagon equation holds exactly for all five-tuples with fourth object
$V$. The run took $39{,}892$ seconds with $2$ workers
(\texttt{deposit/\allowbreak logs/\allowbreak pent\_reduced\_descended.log}).

\emph{By Lemma~\ref{lem:reduction}.} Since every simple object of $\CF$ is a summand of a positive
tensor power of $V$, and the F-matrices are invertible, the pentagon equation for the five-tuples with
fourth object $V$ implies it for all $1{,}030{,}051$ five-tuples. This proves Theorem~\ref{thm:main}(ii).

\emph{Larger exact subsets} (\texttt{pent\_exact\_subset.py}). Independently of the reduction, the
pentagon equation is proved, by the same method and at all $24$ roots, for all $339{,}330$
five-tuples (from all $16$ objects) with $n_{ABCDE}\le50$ (bounds up to $810$ bits, up to $27$ primes,
$38{,}984{,}496$ five-tuple checks, $21{,}721$ seconds with $2$ workers), for all $13{,}137$ five-tuples
of objects from $\{\ob{0000},\ob{0001},\ob{0002},\ob{0010},\ob{1000},\ob{0004},\ob{2000}\}$ (dimensions
up to $395$, bounds up to $391$ bits) and for all $1{,}585$ five-tuples from
$\{\ob{0000},\ob{0001},\ob{1000},\ob{0004},\ob{2000}\}$. No failures. The recursion of
Section~\ref{sec:recursion} uses only pentagons with $V$ in the second or third position, so these
include many pentagons not used to produce the data.

\subsection{Modular data}\label{sec:modular}
From the exact R-symbols and the quantum dimensions~\eqref{eq:qdim}, we computed in $\K$, with
rational arithmetic (\texttt{exact\_checks.py}),
\[
\tilde\theta_X=\frac1{d_X}\sum_Zd_Z\Tr R[X,X,Z],\qquad
\tilde S_{XY}=\sum_Zd_Z\Tr\bigl(R[Y,X,Z]R[X,Y,Z]\bigr),
\]
which in a ribbon category with dimensions $d_X$ are the twists and the unnormalised $S$-matrix.
We found $\tilde\theta_X=q^{(X,X+2\rho)}$ for all $X$, and $\tilde S_{XY}=S_{XY}/S_{\one\one}$ for
all $256$ pairs, where
\[
\frac{S_{XY}}{S_{\one\one}}=\frac{\sum_{w\in W}\varepsilon(w)\,q^{-2(w(X+\rho),Y+\rho)}}
{\sum_{w\in W}\varepsilon(w)\,q^{-2(w\rho,\rho)}}
\]
is the Kac--Peterson formula~\cite{KacPeterson}, $W$ the Weyl group of $F_4$ (of order $1152$).
We also found $R[Y,X,Z]R[X,Y,Z]=\theta_Z\theta_X^{-1}\theta_Y^{-1}\cdot\id$ for all $2{,}302$
triples (the balancing identity). This proves Theorem~\ref{thm:main}(iv). During the construction
the trace of each idempotent $p_X$ was also compared with $d_X$, modulo $p$.
Finally, $\tilde S\tilde S=\mathcal D^2\cdot\id$ holds exactly, with $\mathcal D^2=\sum_Xd_X^2\ne0$
(\texttt{s\_nondegenerate.py}; all objects are self-dual, so the charge conjugation matrix is the
identity); in particular $\tilde S$ is invertible.

\subsection{Modularity}\label{sec:modularity}
We use the data over $\C$ through the embedding $\zeta\mapsto e^{2\pi i/52}$. For a fusion category
over $\C$, the squared norm $|X|^2$ of a simple object and the global dimension
$\dim\cC=\sum_X|X|^2$ are defined in~\cite{ENO}; $\cC$ is pseudo-unitary if $\dim\cC$ equals its
Frobenius--Perron dimension.

\begin{lemma}\label{lem:norm}
Let $X$ be a self-dual simple object with $N_{XX}^\one=1$, in data with identity F-matrices on unit
legs, and put $f=F[X,X,X,X]_{(\one,0,0),(\one,0,0)}$. Then
$(F[X,X,X,X]^{-1})_{(\one,0,0),(\one,0,0)}=f$ and $|X|^2=f^{-2}$.
\end{lemma}

\begin{proof}
Let $a=\phi^{XX\to\one}_0$ and $c$ the vector of $\Hom(\one,X\otimes X)$ with $a\circ c=1$. As in the
proof of Proposition~\ref{prop:rigid}, $(a\otimes\id_X)\circ(\id_X\otimes c)=f\cdot\id_X$; expanding
$\phi^{X\one\to X}_0\circ(\id_X\otimes a)$ in the left-bracketed basis with the inverse F-matrix and
composing with $c\otimes\id_X$ gives $(\id_X\otimes a)\circ(c\otimes\id_X)=f'\cdot\id_X$ with
$f'=(F[X,X,X,X]^{-1})_{(\one,0,0),(\one,0,0)}$. With $c'=c/f$ the first composite is $\id_X$, and the
proof of Proposition~\ref{prop:rigid} shows that the second is then $\id_X$ as well, so $f'=f$. Thus
$(X,a,c')$ exhibits $X$ as its own dual on both sides. With the identity as isomorphism
$X\to X^{**}$, the two traces in the definition of $|X|^2$ in~\cite{ENO} are both $a\circ c'=1/f$, so
$|X|^2=f^{-2}$.
\end{proof}

\begin{proposition}\label{prop:modular}
The braided fusion category defined by the data, with its spherical structure of positive dimensions,
is a modular tensor category with dimensions $d_X$, twists $\theta_X=q^{(X,X+2\rho)}$ and $S$-matrix that
of Kac and Peterson.
\end{proposition}

\begin{proof}
By Theorem~\ref{thm:main}(iii) and Lemma~\ref{lem:norm}, $|X|^2=d_X^2$ for every simple $X$, so
$\dim\cC=\sum_Xd_X^2$. The numbers $d_X$ are positive at our embedding, and
$\sum_ZN_{XY}^Zd_Z=d_Xd_Y$ exactly (\texttt{pseudo\_unitary.py}). So $(d_Z)$ is a positive eigenvector of
the fusion matrix of $V$, which is irreducible because every simple object is a summand of a power of
$V$; by the Perron--Frobenius theorem $d_X$ is the Frobenius--Perron dimension of $X$ for every $X$, and
$\cC$ is pseudo-unitary. By~\cite[Proposition~8.23]{ENO} it has a unique spherical structure with positive
dimensions, and these are the $d_X$. A spherical braided fusion category is ribbon, with twist $\theta$
determined by the braiding and the spherical structure~\cite{BakalovKirillov}; its trace satisfies
$\Tr(\theta_X)=\Tr(c_{X,X})$ and $\Tr(c_{Y,X}c_{X,Y})=\tilde S_{XY}$ for the unnormalised $S$-matrix. For an
endomorphism $u$ of $X\otimes Y$ with matrix $M_Z$ on each $\Hom(X\otimes Y,Z)$ in the basis of fusion
maps, $\Tr(u)=\sum_Zd_Z\Tr M_Z$, and by~\eqref{eq:R} the matrices of $c_{X,X}$ and of $c_{Y,X}c_{X,Y}$ on
$\Hom(X\otimes X,Z)$ and $\Hom(X\otimes Y,Z)$ are $R[X,X,Z]$ and $R[Y,X,Z]R[X,Y,Z]$. So
$\theta_X=\tilde\theta_X$ and the $S$-matrix is $\tilde S$, which equal $q^{(X,X+2\rho)}$ and the
Kac--Peterson values, and $\tilde S^2=\mathcal D^2\cdot\id$, exactly (Section~\ref{sec:modular}). So
$\tilde S$ is invertible and the category is modular.
\end{proof}

As a check of Lemma~\ref{lem:norm} on the data, the identity
\[(F[X,X,X,X]^{-1})_{(\one,0,0),(\one,0,0)}=F[X,X,X,X]_{(\one,0,0),(\one,0,0)}\]
holds for all $16$ objects modulo two primes at all $12$ real roots (\texttt{pseudo\_unitary.py},
\texttt{deposit/logs/pseudo\_unitary.log}).

\subsection{Unitarity}\label{sec:unitarity}
In the conventions of Rowell's survey, where the form is normalised by short roots of length $2$,
$q^2$ is a primitive $\ell$-th root of unity and the alcove is $(\lambda+\rho,\theta)<\ell$, the
category $\cC(\mathfrak g,\ell,q)$ is unitary when the lacing number $m$ divides $\ell$ and
$q=e^{\pi i/\ell}$, by results of Wenzl~\cite{Wenzl} and Xu~\cite{Xu}, as stated
in~\cite[Section~4]{Rowell}. For $\CF$ we have $\ell=26$, $m=2$ and $q=e^{\pi i/26}=\zeta$, so $\CF$ is
unitary at the embedding $\zeta=e^{2\pi i/52}$ used throughout.

\begin{remark}
A unitary structure on skeletal data is a family of positive definite Hermitian forms $G^{XY}_Z$ on the
multiplicity spaces for which every F-matrix is an isometry between the Gram matrices of the two
bases of $\Hom(X\otimes Y\otimes W,T)$ built from them. With the normal form, $G$ can be computed from
the F-symbols with third leg $V$; reconstructed over $\kk$, its entries have at most $47$ bits. The
isometry identity was checked on all $58{,}740$ blocks modulo two (prime, root) pairs, and the forms
are positive definite on all $2{,}271$ multiplicity spaces without a unit leg exactly at the
embeddings $q=\zeta^{\pm1}$ and $q=\zeta^{\pm5}$ of $\kk$, and indefinite on $520$ to $549$ of them at
the other four, with eigenvalues computed numerically at $256$-bit precision
(\texttt{unitary/f4\_forms\_analyse.py}). In a gauge in which all F-matrices
are real orthogonal, every entry whose four multiplicity spaces are one dimensional squares to a gauge
invariant; the square classes of these invariants span a subgroup of $\kk^\times/(\kk^\times)^2$ of rank
$6$, generated by $2,[3],[5],[11],[13],[2][4]$ (detected by quadratic characters at $300$ primes, which
proves their independence; \texttt{unitary/f4\_classes.py}). So any such gauge needs a field of degree at
least $2^6\cdot6=384$ over $\Q$, which is why the data are given in the normal form.
\end{remark}

\section{Autoequivalences and braidings}\label{sec:autoeq}
\subsection{What was known}
Gannon determined the automorphisms of the fusion rings of all affine Lie algebras~\cite{Gannon}.
For $F_4$ at level $4$ he found an automorphism of order $4$, which he calls $\pi_4$, and showed that
$\pi_4,\pi_4^2,\pi_4^3$ are the only nontrivial ones~\cite[Section~3.8 and Theorem~3.F4]{Gannon}; in
our labels
\[
\pi_4=(\ob{0001}\ \ob{1000}\ \ob{2000}\ \ob{0004})(\ob{0010}\ \ob{0100}\ \ob{0020}\ \ob{0003})
(\ob{1001}\ \ob{1002}\ \ob{1010}\ \ob{0012}),
\]
fixing $\ob{0000},\ob{0002},\ob{0011},\ob{0101}$. Edie-Michell computed the monoidal and braided
autoequivalence groups of $\cC(\mathfrak g,k)$ for $\mathfrak g$ of types $A$, $B$, $C$ and
$G$~\cite{EdieMichell}. For $F_4$ he writes (\cite[Section~1]{EdieMichell}, arXiv v2): ``For the
categories $\mathcal C(\mathfrak{so}_{2r},2)$ and $\mathcal C(\mathfrak f_4,4)$ there exist exceptional
auto-equivalences. The author was unable to find ad-hoc methods for constructing these. [\dots{}] The
exceptional auto-equivalence of $\mathcal C(\mathfrak f_4,4)$ remains completely mysterious''
(``exist'' refers to the fusion ring automorphism), and observes in~\cite[Section~7]{EdieMichell} that
$\CF$ is the only other quantum group modular tensor category that could have a reverse braided
autoequivalence. Morrison, Snyder and Thurston~\cite[Section~6]{MorrisonSnyderThurston} list
$(E_8,F_4,\zeta_{52}^9)\cong(F_4,F_4,\zeta_{52}^7)$ among their conjectural isomorphisms of quantum group
categories and remark that it is ``a categorical version of the automorphism of the fusion ring of
$F_4$ level~$4$'' of Gannon, and that composing it with its Galois conjugate gives a braided equivalence
$((F_4)_4,\zeta_{52})\cong((F_4)_4,\zeta_{52}^{25})\cong((F_4)_4,\zeta_{52}^{-1})$ ``and hence a braided
reversing auto-equivalence'', adding that one needs ``to assume the quantum sufficiency conjecture in
order to construct the autoequivalence''. The results below decide these questions for $\CF$
unconditionally, from the data.

\subsection{Notation}
For a fusion ring automorphism $\pi$ we write
\[(\pi\cdot F)[X,Y,W,T]=F[\pi^{-1}X,\pi^{-1}Y,\pi^{-1}W,\pi^{-1}T]\]
and similarly $\pi\cdot R$, with multiplicity indices unchanged; $\sigma_a\cdot(F,R)$ is the data at
$q\mapsto q^a$; and $R^{\mathrm{rev}}[X,Y,Z]=R[Y,X,Z]^{-1}$ is the reverse braiding
$c^{\mathrm{rev}}_{X,Y}=c_{Y,X}^{-1}$. A monoidal equivalence between the categories defined by $F$
and $F'$ that acts on simple objects by $\pi$ is, up to monoidal isomorphism, a unit-preserving gauge
$h$ with $h\cdot(\pi\cdot F)=F'$.

\subsection{Braidings}
\begin{proposition}\label{prop:twobraidings}
Let $(\cC,c)$ be a braided fusion category over $\C$ whose M\"uger centre is trivial, and suppose that
$X\otimes Y$ is not simple for any simple $X,Y\ne\one$. Then the underlying fusion category has exactly
two braidings, $c$ and $c^{\mathrm{rev}}$.
\end{proposition}

\begin{proof}
A braiding $c'$ on $\cC$ defines a fully faithful braided tensor functor $s':\cC\to\mathcal Z(\cC)$,
$X\mapsto(X,c'_{-,X})$, with $s'$ followed by the forgetful functor equal to the identity. Since the
M\"uger centre of $(\cC,c)$ is trivial, the functor $\cC\boxtimes\cC^{\mathrm{rev}}\to\mathcal Z(\cC)$,
$X\boxtimes Y\mapsto(X\otimes Y,\ c_{-,X}\text{ and }c^{\mathrm{rev}}_{-,Y})$, is an equivalence
\cite[Proposition~3.7]{DGNO}. So the simple objects of $\mathcal Z(\cC)$ are the $X\boxtimes Y$, with
underlying object $X\otimes Y$. For simple $Z$ the object $s'(Z)$ is simple, since $s'$ is fully
faithful; so $s'(Z)=X\boxtimes Y$ with $X\otimes Y\cong Z$, hence $s'(Z)\in\{Z\boxtimes\one,
\one\boxtimes Z\}$. If $s'(Z)=Z\boxtimes\one$ and $s'(W)=\one\boxtimes W$ with $Z,W\ne\one$, then
$s'(Z\otimes W)\cong Z\boxtimes W$ is simple, while $Z\otimes W$ has at least two simple summands,
each with nonzero image: a contradiction. So either $s'(Z)\cong Z\boxtimes\one=(Z,c_{-,Z})$ for all
simple $Z$, or $s'(Z)\cong\one\boxtimes Z=(Z,c^{\mathrm{rev}}_{-,Z})$ for all simple $Z$. An isomorphism
$(Z,c'_{-,Z})\cong(Z,c_{-,Z})$ in $\mathcal Z(\cC)$ is a scalar multiple of $\id_Z$ commuting with the
half braidings, so $c'_{-,Z}=c_{-,Z}$; by naturality $c'=c$ in the first case, and $c'=c^{\mathrm{rev}}$
in the second. Finally $c\ne c^{\mathrm{rev}}$, since otherwise every object would be in the M\"uger
centre.
\end{proof}

For $\CF$ and all its Galois conjugates both hypotheses hold, exactly from the data: no tensor product
of two non-unit simple objects is simple, and for every simple $X\ne\one$ some monodromy
$R[Y,X,Z]R[X,Y,Z]$ is a scalar different from $1$ (Section~\ref{sec:modular};
\texttt{autoeq\_checks.py}). Consequently, if $\Phi$ is a monoidal equivalence from $\CF$ to a Galois
conjugate $\sigma_a\cdot\CF$, acting on objects by $\pi$, then $\Phi$ transports $\sigma_a(c)$ to $c$ or
to $c^{\mathrm{rev}}$. The two cases are distinguished by the gauge invariant numbers $R[X,X,Z]$ on the
$100$ one-dimensional spaces, which a braided equivalence preserves.

\subsection{Results}
\begin{theorem}\label{thm:autoeq}
Let $\pi=\pi_4$, and let $H=\{1,5,21,25,27,31,47,51\}\subset(\Z/52)^\times$.
\begin{enumerate}[label=(\roman*)]
\item The fusion ring of $\CF$ has exactly four automorphisms, the powers of $\pi$.
\item $\pi^2$ lifts to a monoidal autoequivalence $\Phi$ of $\CF$, and $\Phi$ transports $c$ to
$c^{\mathrm{rev}}$: it is a braided equivalence $\CF\to\CF^{\mathrm{rev}}$. The automorphisms $\pi$ and
$\pi^3$ do not lift to monoidal autoequivalences.
\item Every monoidal autoequivalence of $\CF$ that fixes all simple objects is monoidally isomorphic to
the identity.
\item Hence the group of monoidal autoequivalences of $\CF$ up to monoidal isomorphism is $\Z/2$,
generated by $\Phi$, and the group of braided autoequivalences is trivial. The underlying fusion
category has exactly two braidings, $c$ and $c^{\mathrm{rev}}$, and they are exchanged by $\Phi$.
\item $\pi$ lifts to a braided equivalence $\CF\to\sigma_5\cdot\CF$, from the category at $q$ to the
category at $q^5$. The data $\sigma_a\cdot(F,R)$ are braided equivalent to $(F,R)$ exactly for $a\in H$,
and for $a\notin H$ the fusion category $\sigma_a\cdot\CF$ is not monoidally equivalent to $\CF$ by any
equivalence. So the $24$ Galois conjugates of $\CF$ form $3$ braided equivalence classes, and the field
of moduli up to relabelling of the objects, for the fusion category and for the braided category, is
the fixed field of $H$, the cubic subfield of $\Q(\zeta_{13})$ (for forms over it see
Section~\ref{sec:forms}).
\end{enumerate}
\end{theorem}

\begin{proof}
(i) Recomputed from the fusion rules by a search over permutations preserving the multiplicities
(\texttt{autoeq\_checks.py}, \texttt{fusion\_auts.py}); it agrees with~\cite[Theorem~3.F4]{Gannon}.

\emph{Monoidal obstructions.} In the blocks $F[V,V,V,V]$ and $F[V,V,V,\ob{0002}]$ all multiplicity spaces
are one dimensional, so every gauge acts on them by $F\mapsto D_1FD_2$ with diagonal $D_i$, and the
numbers $I_{ab}=F_{ab}(F^{-1})_{ba}$ are invariants of the gauge class. Exactly in $\K$, the relation
$I(\pi^j\cdot F)=\sigma_a(I(F))$ holds for $a\in\{1,25,27,51\}$ when $j=0,2$, and only for
$a\in\{5,21,31,47\}$ when $j=1,3$ (\texttt{nf\_witness.py}). So $\pi\cdot F$ and $\pi^3\cdot F$ are not gauge
equivalent to $F$, which is the second half of (ii), and $\sigma_a\cdot F$ is not gauge equivalent to
$\pi^j\cdot F$ for any $j$ when $a\notin H$.

\emph{Monoidal lifts.} Since $F$ has entries in $\kk$, $\sigma_a\cdot F=F$ for $a\in\{1,25,27,51\}$.
Gauges $h$ with $h\cdot(\pi^2\cdot F)=F$ and $h\cdot(\pi^3\cdot F)=\sigma_5\cdot F$ on all blocks with
third leg $V$ are computed modulo primes from the normal form (with $g$ the normal-form gauge of
$\pi^j\cdot F$, $h=g^{-1}$, suitably conjugated), reconstructed over $\Lf$, and the identities proved
exactly on all $3{,}451$ such blocks with a height bound, with every $h^{XY}_Z$ invertible
(\texttt{f4\_relabel\_exact.py} on this table, $537$ and $530$ s; \texttt{deposit/logs/nf\_relabel\_exact.log}).
Both sides satisfy the pentagon equation, so by Lemma~\ref{lem:Hdetermines} the identities hold on the whole
table: $\pi^2\cdot F\sim F$ and $\pi^3\cdot F\sim\sigma_5\cdot F$. Applying $\pi$ and then
$\sigma_{21}=\sigma_5^{-1}$ to the second gives
$\pi\cdot F\sim\sigma_{21}\cdot F=\sigma_5\cdot(\sigma_{25}\cdot F)=\sigma_5\cdot F$. Together with the
cases $a\in\{1,25,27,51\}$, this gives a monoidal equivalence $\sigma_a\cdot\CF\simeq\CF$ for every
$a\in H$.

\emph{Braidings.} For each such equivalence, Proposition~\ref{prop:twobraidings} leaves two possibilities
for the transported braiding, and the numbers $R[X,X,Z]$ on one-dimensional spaces decide between
them, exactly (\texttt{autoeq\_checks.py}): $R[\pi^jX,\pi^jX,\pi^jZ]=\sigma_a(R[X,X,Z])$ for all such
$X,Z$ holds for $(j,a)$ with $j=0$, $a\in\{1,27\}$; $j=2$, $a\in\{25,51\}$; $j=1$, $a\in\{21,47\}$;
$j=3$, $a\in\{5,31\}$; and $R[\pi^jX,\pi^jX,\pi^jZ]=\sigma_a(R[X,X,Z])^{-1}$ holds for the complementary
pairs ($j=0$, $a\in\{25,51\}$; $j=2$, $a\in\{1,27\}$; and so on), and for no other pairs. With $j=2$, $a=1$
this is the first half of (ii): $\Phi$ transports $c$ to $c^{\mathrm{rev}}$. With $j=1$, $a=21$ it says
$\sigma_5(R[\pi X,\pi X,\pi Z])=R[X,X,Z]$, so the monoidal equivalence $\CF\to\sigma_5\cdot\CF$ acting by
$\pi$ preserves the braiding; this is the first statement of (v). For $a\in H$, an equivalence $\sigma_a\cdot\CF\simeq\CF$
that reverses the braiding can be composed with $\Phi$, so in all cases $\sigma_a\cdot(F,R)$ is braided
equivalent to $(F,R)$.

(iii) The normal form of Section~\ref{sec:nfrule} fixes a gauge $g_F$ with $g_F\cdot F=\mathrm{NF}(F)$
(here $g_F=1$, since $F$ is in normal form) by conditions on $g\cdot F$, and the set of gauges satisfying these conditions is the coset $Tg_F$ of the
torus $T=\{g^{XY}_Z=u_Xu_Y/u_Z\}$: in step (2) the integer system has $89$ variables, rank $74$ and $74$
unit pivots, so its solutions form a coset of a connected torus of dimension $15$, which is the
restriction of $T$ ($T$ has dimension $15$ and fixes all F-symbols); steps (3) and (4) determine the gauge
on every one of the remaining $2{,}182$ multiplicity spaces without a unit leg uniquely from the earlier
ones (\texttt{stab\_check.py}). If $h\cdot F=F$ then $g_Fh$ satisfies the conditions, so $g_Fh\in Tg_F$, and
since $T$ consists of scalars on each multiplicity space it is central in the gauge group, so $h\in T$.
A gauge in $T$ is the monoidal structure of the identity functor transported by the natural isomorphism
$(u_X)$, hence monoidally isomorphic to the identity.

(iv) By (iii) the map from monoidal autoequivalences up to isomorphism to fusion ring automorphisms is
injective; by (ii) its image is $\{1,\pi^2\}$. A braided autoequivalence fixes the numbers
$R[X,X,Z]$, so it is not $\Phi$; by (iii) it is the identity. The statement about braidings is
Proposition~\ref{prop:twobraidings}.

(v) The braided equivalence for $\pi$ is shown above. For $a\notin H$, the monoidal obstruction above
excludes every relabelling, and a monoidal equivalence acts on objects by a fusion ring automorphism.
The subgroup $H$ has order $8$ and index $3$ in $(\Z/52)^\times$, and its fixed field is the cubic subfield
of $\Q(\zeta_{13})$.
\end{proof}

As an independent confirmation of the braided statements, the canonical R-symbols
$g^{YX}_Z\,R[X,Y,Z]\,(g^{XY}_Z)^{-1}$ of $(\pi^j\cdot F,\pi^j\cdot R)$ and of $(F,R)$ or $(F,R^{\mathrm{rev}})$
at the conjugate roots were compared directly, modulo three primes at all $24$ roots, with the results
above (\texttt{braid\_transport\_f4.py}).

In particular, the equivalences predicted in~\cite[Section~6]{MorrisonSnyderThurston}, assuming the
quantum sufficiency conjecture, exist for $\CF$: the braided equivalence
$((F_4)_4,\zeta_{52})\cong((F_4)_4,\zeta_{52}^{5})$ acting by Gannon's $\pi_4$ (Theorem~\ref{thm:autoeq}(v)), and a braid-reversing
autoequivalence, $\Phi$. By Theorem~\ref{thm:autoeq}(iv), $\Phi$ is, up to monoidal isomorphism, the only
nontrivial monoidal autoequivalence.

\subsection{Forms over the cubic field}\label{sec:forms}
Let $E\subset\Kt$ be the cubic subfield, $E=\Q(\vartheta)$ with
$\vartheta=\zeta_{13}+\zeta_{13}^5+\zeta_{13}^8+\zeta_{13}^{12}$ and
$\vartheta^3+\vartheta^2-4\vartheta+1=0$, and let $s\in\mathrm{Gal}(\Kt/E)$ be
$\zeta_{13}\mapsto\zeta_{13}^5$, the restriction of $\sigma_5$. Then $s$ generates
$\mathrm{Gal}(\Kt/E)\cong\Z/4$, $s^2$ is complex conjugation, and $s|_\kk$ generates
$\mathrm{Gal}(\kk/E)\cong\Z/2$. Let $\mathcal A$ be the fusion category over $\kk$ defined by $F$, with the
labels of Table~\ref{tab:objects}, and $\mathcal A_M=\mathcal A\otimes_\kk M$ for fields $M\supseteq\kk$.
Following Etingof and Gelaki~\cite[Definition~3.1]{EtingofGelakiDescent}, a \emph{form} of $\mathcal A_M$
over a subfield $M_0\subseteq M$ is a fusion category $\mathcal D$ over $M_0$ with an $M$-linear monoidal
equivalence $\mathcal D\otimes_{M_0}M\simeq\mathcal A_M$, where $\otimes_{M_0}M$ is extension of scalars
followed by idempotent completion; it is \emph{split} if every simple object of $\mathcal D$ has
endomorphism algebra $M_0$. Braided forms are defined in the same way. The fusion category $\mathcal A$ is
a split form over $\kk$, and by Section~\ref{sec:fields} no smaller field carries a form with the labels
fixed; $E$ is the field of moduli up to relabelling (Theorem~\ref{thm:autoeq}(v)).

For $\gamma\in\mathrm{Aut}(M)$, a $\gamma$-semilinear monoidal autoequivalence of $\mathcal A_M$ is the same
as a monoidal equivalence $\gamma\cdot\mathcal A_M\to\mathcal A_M$~\cite[Remark~3.2]{EtingofGelakiDescent}; if
it acts on simple objects by $p$, it is, in the notation above, a gauge $h$ with
$h\cdot(p\cdot(\gamma\cdot F))=F$. In this direction the braided equivalence $\CF\to\sigma_5\cdot\CF$ of
Theorem~\ref{thm:autoeq}(v), which acts by $\pi$, is the inverse of a $\sigma_5$-semilinear braided
autoequivalence acting by $\pi^{-1}=\pi^3$.

\begin{lemma}\label{lem:twisted}
Let $M\supseteq\kk$ and $\gamma\in\mathrm{Aut}(M)$ with $\gamma(\kk)=\kk$.
\begin{enumerate}[label=(\alph*),leftmargin=2em]
\item A $\gamma$-semilinear monoidal autoequivalence of $\mathcal A_M$ acting on simple objects by $p$
exists if and only if $\gamma|_\kk=\id$ and $p\in\{1,\pi^2\}$, or $\gamma|_\kk=s|_\kk$ and
$p\in\{\pi,\pi^3\}$.
\item Two of them acting by the same $p$ are isomorphic by a unique monoidal isomorphism.
\end{enumerate}
\end{lemma}

\begin{proof}
A semilinear equivalence preserves the dimensions of the spaces $\Hom(X\otimes Y,Z)$, so $p$ is a
fusion ring automorphism, $p=\pi^j$. (b) If $\Psi,\Psi'$ act by $p$, then $\Psi^{-1}\Psi'$ is $M$-linear
and fixes the simple objects, so it is a gauge $g$ with $g\cdot F=F$. The proof of
Theorem~\ref{thm:autoeq}(iii) uses only field operations and applies over $M$: $g$ lies in the torus $T$,
so $\Psi'\cong\Psi$. A monoidal automorphism of the identity functor is a family of scalars $u_X$ with
$u_Xu_Y=u_Z$ whenever $N_{XY}^Z\ne0$; since $N_{VV}^V=1$, $u_V=1$, and since every simple object is a
summand of a tensor power of $V$, $u=1$. So the isomorphism is unique.
(a) If $\gamma|_\kk=\id$ then $\gamma\cdot F=F$, and the statement is Theorem~\ref{thm:autoeq}(ii). If
$\gamma|_\kk=s|_\kk$, then $\gamma\cdot F=\sigma_5\cdot F$, and $\sigma_5\cdot F$ is gauge equivalent to
$\pi\cdot F$ and to $\pi^3\cdot F$ (Theorem~\ref{thm:autoeq}; both gauges are also computed and proved
exactly on all blocks with third leg $V$ by \texttt{relabel\_exact\_nf.py}, \texttt{deposit/logs/forms\_relabel\_pi1.log}
and \texttt{deposit/logs/forms\_relabel\_pi3.log}, hence on the whole table by Lemma~\ref{lem:Hdetermines}); so
$p=\pi^3$ and $p=\pi$ occur. If $p\in\{1,\pi^2\}$ occurred, then $\gamma\cdot F\sim\pi^{-j}\cdot F\sim F$,
so $\gamma$ would fix $\mathrm{NF}(F)=F$, whose entries generate $\kk$. Every other $\gamma|_\kk$ is the
restriction of some $\sigma_a$ with $a\notin H$, and $\sigma_a\cdot F$ is not gauge equivalent to any
$\pi^j\cdot F$ (proof of Theorem~\ref{thm:autoeq}).
\end{proof}

So for $M/E$ finite Galois with $M\supseteq\kk$, and $\varepsilon\colon\mathrm{Gal}(M/E)\to\mathrm{Gal}(\kk/E)=\Z/2$
the restriction, the isomorphism classes of semilinear monoidal autoequivalences of $\mathcal A_M$ over $E$
form the group $G_M=\{(\gamma,j)\in\mathrm{Gal}(M/E)\times\Z/4:j\equiv\varepsilon(\gamma)\bmod2\}$, and the
monoidal automorphisms of the identity functor are trivial. Let $\Psi_j$ ($j=1,3$) be the $s$-semilinear
monoidal autoequivalence of $\mathcal A_{\Kt}$ acting by $\pi^j$, with gauge entries in $\kk$. Then
$\Psi_j^4$ is $\Kt$-linear and fixes the simple objects, so by Lemma~\ref{lem:twisted}(b) there is a unique
monoidal isomorphism $\Psi_j^4\Rightarrow\id$, given by scalars $u_X=u^{(j)}_X$ with $u_\one=1$. The program
\texttt{descent\_u.py} computes the gauges of $\Psi_j$ and $\Psi_j^4$ exactly, checks that the gauge of
$\Psi_j^4$ is $u_Xu_Y/u_Z$ on all $1{,}394$ one-dimensional multiplicity spaces (and, modulo a prime, on all
$2{,}302$), and finds $u_{\pi^jX}=s(u_X)$ and $u_X\in E$ for the four $\pi$-fixed $X$. For $j=3$,
\begin{gather*}
u_{\ob{0002}}=\tfrac1{169}(244\vartheta^2+583\vartheta-178),\qquad
u_{\ob{0101}}=\tfrac1{169}(-14876486\vartheta^2-35364425\vartheta+10801957),\\
u_{\ob{0011}}=\tfrac1{28561}(-88526336\vartheta^2-210445056\vartheta+64279808).
\end{gather*}

\begin{proposition}\label{prop:forms}
Let $M/E$ be finite Galois with $M\supseteq\kk$.
\begin{enumerate}[label=(\alph*),leftmargin=2em]
\item Forms of $\mathcal A_M$ over $E$, up to equivalence, correspond to homomorphisms
$\chi\colon\mathrm{Gal}(M/E)\to\Z/4$ with $\chi\equiv\varepsilon\bmod2$; in the form $\mathcal D_\chi$ the
Galois group permutes the simple objects of $\mathcal A_M$ by $\gamma\mapsto\pi^{\chi(\gamma)}$. In
particular $\mathcal A$ has no form over $E$, and no form $\mathcal D$ over $E$ of any $\mathcal A_M$ has
$\mathcal D\otimes_E\kk\simeq\mathcal A$. For $M=\Kt$ there are exactly two forms, $\mathcal D_1$ and
$\mathcal D_3$ ($\chi_j(s)=j$, descent data $\Psi_j$), and they are not equivalent.
\item $\mathcal D_1$ and $\mathcal D_3$ have seven simple objects: the unit, with endomorphism algebra $E$;
one for each $\pi$-orbit of size four, with endomorphism algebra $\Kt$ and base change the sum of the
orbit; and for each $X\in\{\ob{0002},\ob{0011},\ob{0101}\}$ an object $S_X$ with
$S_X\otimes_E\Kt\cong X\oplus X$ whose endomorphism algebra is the quaternion division algebra $B_X$ over
$E$ ramified exactly at the two real places at which $u_X<0$: at $\vartheta\approx-2.651$ and $0.274$ for
$\ob{0002}$, at $-2.651$ and $1.377$ for $\ob{0011}$, at $0.274$ and $1.377$ for $\ob{0101}$ (for both $j$).
No form of $\mathcal A_{\Kt}$ over $E$ is split.
\item Forms of $\mathcal A_M$ over $\kk$ are $\mathcal A$ and, for each quadratic extension
$L=\kk(\sqrt d)\subseteq M$, one form $\mathcal A^{(L)}$, on which $\mathrm{Gal}(L/\kk)$ acts through $\Phi$. Its
simple objects are the unit, one object with endomorphism algebra $L$ for each of the six
$\pi^2$-orbits of size two, and one for each $X\in\{\ob{0002},\ob{0011},\ob{0101}\}$, with endomorphism
algebra the division algebra in the class of the quaternion algebra $(d,u_X)_\kk$. For $L=\Kt$,
$\mathcal A^{(\Kt)}=\mathcal D_j\otimes_E\kk$, and these are $B_X\otimes_E\kk$, ramified at four real places.
\item For a number field $M_0$, $\mathcal A\otimes_\kk\overline\Q$ has a form over $M_0$ if and only if
$M_0\supseteq E$, and a split form if and only if $M_0\supseteq\kk$.
\end{enumerate}
\end{proposition}

\begin{proof}
(a) By~\cite[Sections~3.2--3.5, 3.7, Propositions~3.5 and~3.8]{EtingofGelakiDescent}, forms over $E$
correspond to splittings $\gamma\mapsto(\gamma,\chi(\gamma))$ of $G_M\to\mathrm{Gal}(M/E)$ together with
coherence data whose obstruction lies in $H^3$ and whose choices form a torsor under $H^2$ of
$\mathrm{Gal}(M/E)$ with coefficients in the monoidal automorphisms of the identity functor; both groups
vanish. Unframed classes are orbits of $\mathrm{Aut}_\otimes(\mathcal A_M)=\{1,\Phi\}$, which is central in
$G_M$, so they agree with framed classes. A form $\mathcal D$ with $\mathcal D\otimes_E\kk\simeq\mathcal A$ would
have $\chi=0$ on $\mathrm{Gal}(M/\kk)$; but for $\varepsilon(\gamma)=1$, $\gamma^2\in\mathrm{Gal}(M/\kk)$ and
$\chi(\gamma^2)=2\chi(\gamma)=2$ (compare~\cite[Example~3.4]{EtingofGelakiDescent}). An equivalence
$\mathcal D_1\simeq\mathcal D_3$ would give $a\in\{1,\Phi\}$ with $a\Psi_1a^{-1}\cong\Psi_3$, impossible as
$\Phi$ is central and $\Psi_1,\Psi_3$ act by $\pi,\pi^3$.

(b) The form is the category of objects $Y$ of $\mathcal A_{\Kt}$ with $\alpha\colon\Psi_j(Y)\to Y$ satisfying
$\alpha\circ\Psi_j(\alpha)\circ\Psi_j^2(\alpha)\circ\Psi_j^3(\alpha)=u$ on $Y$, with the invariant morphisms
\cite[proof of Proposition~3.5]{EtingofGelakiDescent}. An orbit of size four gives a simple object with
endomorphism algebra $\Kt$. For fixed $X$, such an $\alpha$ on $X^{\oplus n}$ is $\alpha\in GL_n(\Kt)$, and
$v\mapsto\alpha\,s(v)$ makes $\Kt^n$ a module over the cyclic algebra $\Lambda_X=(\Kt/E,s,u_X)$ of degree $4$,
whose module endomorphisms are the morphisms of the form. (With the opposite convention $u_X^{-1}$ replaces
$u_X$, which gives the same class below.) At a real place of $E$ the extension $\Kt$ is complex, so the
local invariant is $\tfrac12$ where $u_X<0$ and $0$ where $u_X>0$. The norms of the $u_X$ are divisible only
by $2$ and $13$; $2$ is inert in $\Kt$, with $u_X$ of valuation $0$, $4$ or $8$ there, so the local
invariant at $2$ is $0$; by reciprocity it is $0$ at the prime above $13$ as well. So $[\Lambda_X]$ has order
$2$, hence index $2$ (Albert--Brauer--Hasse--Noether), $\Lambda_X\cong M_2(B_X)$, and the simple object is
$X^{\oplus2}$ with endomorphism algebra $B_X^{\mathrm{op}}\cong B_X$. For the unit $u_\one=1$ and $\Lambda_\one$ is split.

(c) Over $\kk$ the datum sends $\gamma\in\mathrm{Gal}(M/\kk)$ nontrivial on $L$ to the $\gamma$-semilinear
autoequivalence acting by $\pi^2$ with the gauge of $\Psi_j^2$ (entries in $\kk$); its square has the same
scalars $u_X$, and the same argument applies with $(L/\kk,\tau,u_X)=(d,u_X)_\kk$. For $L=\Kt$,
$\chi_j(s^2)=2$, so $\mathcal A^{(\Kt)}=\mathcal D_j\otimes_E\kk$; $d=(\zeta_{13}-\zeta_{13}^{-1})^2$ is
totally negative, so $(d,u_X)_\kk$ is ramified at the four real places over those of $B_X$.

(d) If $M_0\not\supseteq E$, some $\gamma$ fixing $M_0$ restricts on $\kk$ to $\sigma_a$, $a\notin H$, and by
Lemma~\ref{lem:twisted} there is no semilinear autoequivalence over $\gamma$, so no
form~\cite[Section~3.2]{EtingofGelakiDescent}. If $M_0\supseteq E$, take $\chi_1$ composed with the restriction
$\mathrm{Gal}(M/M_0)\to\mathrm{Gal}(\Kt/E)$, for $M\supseteq M_0\Kt$ finite Galois over $\Q$. In a split form every simple object of $\mathcal A_M$ is
fixed, so $\chi=0$ and $\varepsilon=0$ on $\mathrm{Gal}(M/M_0)$, that is $M_0\supseteq\kk$.
\end{proof}

The signs, Hasse invariants and indices are computed in PARI/GP by \texttt{brauer.gp} (written by
\texttt{brauer.py}); its \texttt{alghassei}, \texttt{alghassef} and \texttt{algindex} agree with the argument
above, and the signs, norms and the valuations at $2$ were also recomputed independently from the values of
$u_X$.

\begin{proposition}\label{prop:braidedforms}
Let $c$ be the braiding of $\mathcal A_{\Kt}$ given by $R$.
\begin{enumerate}[label=(\alph*),leftmargin=2em]
\item For $\gamma\in\mathrm{Gal}(\overline\Q/E)$ there is exactly one braided $\gamma$-semilinear
autoequivalence of $(\mathcal A_{\overline\Q},c)$ up to isomorphism; it acts by $1$, $\pi^3$, $\pi^2$, $\pi$
according as $\gamma|_{\Kt}=1,s,s^2,s^3$. In particular $\Psi_3$ is braided and $\Psi_1$ transports $c$ to
$c^{\mathrm{rev}}$.
\item For a number field $M_0$, $(\mathcal A_{\Kt},c)$ has a braided form over $M_0$ if and only if
$M_0\supseteq E$, and it is unique up to braided equivalence. Over $E$ it is $\mathcal D_3$ with the
descended braiding, and $\mathcal D_1$ admits no braiding; over $\kk$ it is $\mathcal A^{(\Kt)}$. A split braided
form exists over $M_0$ if and only if $M_0\supseteq\Kt$.
\end{enumerate}
\end{proposition}

\begin{proof}
(a) A $\gamma$-semilinear monoidal autoequivalence acting by $\pi^j$ transports $c$ to $c$ or to
$c^{\mathrm{rev}}$ (Proposition~\ref{prop:twobraidings}), and it is braided if and only if
$R[\pi^jX,\pi^jX,\pi^jZ]=\gamma(R[X,X,Z])$ on the $100$ one-dimensional spaces. With $\gamma|_{\Kt}=\sigma_a|_{\Kt}$
this is the table in the proof of Theorem~\ref{thm:autoeq}: $(j,a)$ with $a\equiv1,5,12,8\pmod{13}$
gives $j=0,3,2,1$ (also checked at all $24$ roots modulo a prime by \texttt{descent\_u.py}). Uniqueness
follows from Lemma~\ref{lem:twisted}(b). (b) By~\cite[Section~3.9]{EtingofGelakiDescent} braided forms are
classified in the same way with braided autoequivalences, whose classes over $E$ map isomorphically onto
the Galois group by (a); so there is exactly one splitting, $\chi(s)=3$. A braiding on a form is a
braiding of $\mathcal A_{\Kt}$ invariant under the descent datum; $c$ and $c^{\mathrm{rev}}$ are invariant
under $\Psi_3$ and exchanged by $\Psi_1$. A split braided form fixes the labels, so $\gamma|_{\Kt}=1$ for all
$\gamma$ fixing $M_0$.
\end{proof}

So the three braided equivalence classes of Galois conjugates in Theorem~\ref{thm:autoeq}(v) are the base
changes of the single braided fusion category $\mathcal D_3$ over $E$ along the three embeddings $E\to\C$.
The field of moduli up to relabelling, $E$, is a field of definition, but only of non-split forms; split
forms need $\kk$ for the fusion category and $\Kt$ for the braided category.

All forms of the fusion category over $E$ are described as follows. For $e\in E^\times$ let $\psi_e$ be
the quadratic character of $E(\sqrt e)$, $\chi_e=\chi_1+2\psi_e$, $\mathcal D_e$ the form of
$\mathcal A_{\overline\Q}$ over $E$ with character $\chi_e$ (Proposition~\ref{prop:forms}(a)), and $M_e$ the
fixed field of $\ker\chi_e$. Here $u_X=u^{(1)}_X$; the ratios $u^{(1)}_X/u^{(3)}_X$ are squares in $E$
(\texttt{endo.gp}), so $u^{(3)}_X$ gives the same quaternion algebras.

\begin{proposition}\label{prop:endo}
\begin{enumerate}[label=(\alph*),leftmargin=2em]
\item $e\mapsto\mathcal D_e$ is a bijection from $E^\times/E^{\times2}$ to the forms of $\mathcal A_{\overline\Q}$ over
$E$, with $\mathcal D_{13}=\mathcal D_3$. $M_e$ is a cyclic quartic extension of $E$ containing $\kk$, and
$\mathcal D_e\otimes_EF$ is split if and only if $F\supseteq M_e$. $M_e=\Kt$ exactly for $e\in\{1,13\}$, and
these are the only $e$ for which $M_e$ is unramified outside $13$ and $\infty$; $M_{-1}=\Q(\zeta_{52})^+$.
\item $\mathcal D_e$ has seven simple objects: the unit, with endomorphism algebra $E$; one object for each
$\pi$-orbit of size four, with endomorphism algebra $M_e$; and for each
$X\in\{\ob{0002},\ob{0011},\ob{0101}\}$ an object $S_X$ whose endomorphism algebra is the division algebra
in the class of the quaternion algebra $(-e,u_X)_E$ (with $S_X\otimes_E\overline\Q$ equal to $X\oplus X$ or
$X$). The local invariant of $(-e,u_X)_E$ is $\tfrac12$ at a real place $v$ exactly when $e(v)>0$ and
$u_X(v)<0$; it is $0$ at the prime above $13$; at a prime $\mathfrak q\nmid26$ it is $\tfrac12$ exactly when
$v_{\mathfrak q}(e)$ is odd and $u_X$ is not a square modulo $\mathfrak q$; at $2$ it is fixed by reciprocity.
\item $[(-e,u_{\ob{0101}})_E]=[(-e,u_{\ob{0002}})_E]+[(-e,u_{\ob{0011}})_E]$ in $\mathrm{Br}(E)$, so the number of
$X$ with $\End(S_X)=E$ is $0$, $1$ or $3$, and all three cases occur. It is $3$ if and only if $e$ is totally
negative and every $u_X$ is a square modulo every $\mathfrak q\nmid26$ at which $v_{\mathfrak q}(e)$ is odd.
For example $\mathcal D_{-1}$, split by $\Q(\zeta_{52})^+$, has only commutative endomorphism algebras:
$E$ for four simple objects and $\Q(\zeta_{52})^+$ for three.
\item $\mathcal D_e$ admits a braiding if and only if $e=13$.
\end{enumerate}
\end{proposition}

\begin{proof}
(a) Every lift of $\varepsilon$ is $\chi_1+2\psi$, and Kummer theory identifies the $\psi$ with
$E^\times/E^{\times2}$; $\chi_3=3\chi_1=\chi_1+2\psi_{13}$, since $\kk=E(\sqrt{13})$. Since $\chi_e$ is odd off
$\mathrm{Gal}(\overline\Q/\kk)$, it is surjective with kernel inside $\mathrm{Gal}(\overline\Q/\kk)$. Over
$F\supseteq M_e$ the character vanishes and the form becomes $\mathcal A_F$; conversely the orbit objects
need $M_e\subseteq F$. At a prime $\mathfrak q\nmid13$, $\chi_1$ is unramified, so $\chi_e$ is ramified exactly
when $\psi_e$ is. As $E$ has class number one, $\psi_e$ is unramified outside $\{2,13,\infty\}$ exactly for
$e$ in the group generated by $-1$, the fundamental units, $2$ and a generator of the prime above $13$, and
among these only $1$ and $13$ give $\psi_e$ unramified at $2$ ($2$ is inert in $E$, so an odd power of $2$
ramifies; for the others, \texttt{endo.gp}). At a real place
$\chi_e(c_v)=2+2\psi_e(c_v)$, which vanishes for $e=-1$, so $M_{-1}$ is the real cyclic quartic extension of $E$
in $\Kt(i)=\Q(\zeta_{52})$.

(b) For $e\notin\{1,13\}$ take $M=\Kt(\sqrt e)$, with $\mathrm{Gal}(M/E)=\langle\tilde s\rangle\times\langle r\rangle$,
$\tilde s|_{\Kt}=s$, $\tilde s(\sqrt e)=\sqrt e$, $r$ fixing $\Kt$. The datum $\tilde s\mapsto\Psi_1$,
$r\mapsto$ the $r$-semilinear autoequivalence acting by $\pi^2$ with the gauge of $\Psi_1^2$ (entries in
$\kk$) has commuting generators, and both $\Psi_1^4$ and the square of the second act on a fixed $X$ by
$u_X$. So the algebra of the proof of Proposition~\ref{prop:forms}(b) becomes
$(\Kt/E,s,u_X)\otimes_E(e,u_X)_E$, of class $[(-1,u_X)_E]+[(e,u_X)_E]=[(-e,u_X)_E]$, because $B_X$ and
$(-1,u_X)_E$ have the same local invariants. The local rule is that of the Hilbert symbol: $u_X$ is a unit
away from $2$ and $13$, and a square in the completion at the prime above $13$ (\texttt{endo.gp}).

(c) $N_{\ob{0002}\,\ob{0011}}^{\ob{0101}}=3$. The $\kk$-linear autoequivalence acting by $\pi^2$ with the gauge
of $\Psi_1^2$ acts on this space by a matrix $g$ with $g^2=(u_{\ob{0002}}u_{\ob{0011}}/u_{\ob{0101}})\,I_3$, so
this ratio is a square in $\kk$ (its cube is $\det(g)^2$); it is a square in $E$ (\texttt{endo.gp}), which gives
the relation. Every real place has some $u_X<0$, which gives the criterion; the three cases occur for
$e=1$, $e=-\vartheta^2-\vartheta+3$ (a fundamental unit) and $e=-1$ (\texttt{endo.gp}).

(d) A braiding on $\mathcal D_e$ is $c$ or $c^{\mathrm{rev}}$ invariant under the descent datum, that is every
$\gamma$ acts by the braided lift, $\chi_e=\chi_3$ (Proposition~\ref{prop:braidedforms}(a)).
\end{proof}

So among all forms over $E$ two properties single out $\Kt$ and $\mathcal D_3$: $\Kt$ is the only splitting field
$M_e$ unramified outside $13$ and $\infty$, and $\mathcal D_3$ is the only form with a braiding; the
quaternion endomorphism algebras are forced for the braided form, but not for the fusion category. The
twisted forms over $\kk$ behave in the same way, with $(d,u_X)_\kk$ in place of $(-e,u_X)_E$; an enumeration
of the totally real quadratic extensions of $\kk$ of small conductor, and the split patterns for all
$\{2,13,\infty\}$-units $e$, are in \texttt{f4-forms/} (\texttt{kforms.gp}, \texttt{endo.gp}).

\emph{Real forms.} Let $\iota_b$, $1\le b\le6$, be the real place of $\kk$ with
$\zeta_{13}+\zeta_{13}^{-1}\mapsto2\cos(2\pi b/13)$; it gives the data at $q=\zeta^w$ for $w\equiv\pm b\pmod{13}$.
Over $\mathbb R$ the category $\mathcal A\otimes_{\iota_b}\C$ has two forms, classified by
$H^1(\mathrm{Gal}(\C/\mathbb R),\{1,\Phi\})=\Z/2$: the split form $\mathcal A\otimes_{\iota_b}\mathbb R$, and the form
$\mathcal A^{(\Kt)}\otimes_{\iota_b}\mathbb R$ in which complex conjugation acts by $\Phi$. The latter is braided
(complex conjugation composed with $\Phi$ is the pair $(j,a)=(2,51)$ in the proof of
Theorem~\ref{thm:autoeq}); its unit has endomorphism algebra $\mathbb R$, the six objects over the
$\pi^2$-pairs have $\C$, and a fixed $X$ has $\mathbb H$ if $u_X(\iota_b)<0$ and $\mathbb R$ otherwise:
\[
\begin{array}{c|c|ccc}
b & \text{place of }E & \ob{0002} & \ob{0011} & \ob{0101}\\\hline
1,5 & \vartheta\approx0.274 & \mathbb H & \mathbb R & \mathbb H\\
2,3 & \vartheta\approx1.377 & \mathbb R & \mathbb H & \mathbb H\\
4,6 & \vartheta\approx-2.651 & \mathbb H & \mathbb H & \mathbb R
\end{array}
\]
The places $b=1,5$ are those at which $\CF$ is unitary (Section~\ref{sec:unitarity}); there the antilinear braided
autoequivalence (complex conjugation)$\circ\Phi$ squares to a negative multiple of the identity on
$\ob{0002}$ and $\ob{0101}$ and to a positive one on $\ob{0011}$.

\section{Non-integrality at 13}\label{sec:integrality}
In $\K=\Q(i,\zeta_{13})$ the prime $13$ factors as $13=\mathfrak P_1^{12}\mathfrak P_2^{12}$
($13$ splits in $\Q(i)$ and is totally ramified in $\Q(\zeta_{13})$). Let $v$ denote either of the
two valuations, normalised by $v(13)=1$, extended to any finite extension of $\K$.

\begin{lemma}\label{lem:qint}
For $0<n<26$, $v([n])=1/12$ if $n$ is even and $v([n])=0$ if $n$ is odd. Consequently
\[
v(d_\lambda)=\tfrac1{12}\Bigl(\#\{\alpha>0:(\lambda+\rho,\alpha)\text{ even}\}-\#\{\alpha>0:(\rho,\alpha)\text{ even}\}\Bigr).
\]
\end{lemma}

\begin{proof}
$q^n-q^{-n}=q^{-n}(\zeta_{26}^n-1)$ with $\zeta_{26}=q^2$. For $0<n<26$ the root of unity
$\zeta_{26}^n$ has order $13$ if $n$ is even and order $26$ or $2$ if $n$ is odd. Now
$1-\xi$ has valuation $1/12$ for $\xi$ of order $13$, and is a unit at $13$ for $\xi$ of order $26$
($1-\xi=1+\xi'$ with $\xi'$ of order $13$, and $\prod(1+\xi')=\Phi_{13}(-1)\cdot(-1)^{12}=1$) or
$2$. The denominator $q-q^{-1}$ is a unit. For the dimension formula note that all pairings
$(\lambda+\rho,\alpha)$ for $\lambda$ in the alcove lie in $(0,26)$.
\end{proof}

For $\lambda=\ob{0011}$ all $24$ pairings $(\lambda+\rho,\alpha)$ are even, while $16$ of the
$(\rho,\alpha)$ are, so $v(d_{\ob{0011}})=8/12=2/3$; for the other $15$ simple objects $v(d_\lambda)=0$
(Table~\ref{tab:objects}). Independently, the $13$-adic valuation of $N_{\K/\Q}(d_{\ob{0011}})$ is
$16=12\cdot(2/3+2/3)$, and that of the other norms is $0$ (\texttt{exact\_checks.py}). Explicitly,
in the basis $1,c_1,\dots,c_{11}$ of $\Lf$, $c_k=\zeta^k+\zeta^{-k}$,
\[
d_{\ob{0011}}=7+14c_2-4c_4+4c_6+14c_{10},\qquad
\frac1{d_{\ob{0011}}}=\frac{2c_2+2c_{10}-5}{13}.
\]

\begin{proposition}\label{prop:integrality}
Let $\tilde F$ be obtained from the F-symbols of Theorem~\ref{thm:main} by any gauge
transformation, with the $g^{XY}_Z$ invertible matrices over any finite extension $L$ of $\K$. Then
for each prime of $L$ above $13$ some F-symbol of $\tilde F$ is not integral at that prime.
\end{proposition}

\begin{proof}
Let $X=\ob{0011}$, so $N_{XX}^\one=1$ and $N_{X\one}^X=N_{\one X}^X=1$. Write $a=g^{XX}_\one$,
$b=g^{\one X}_X$, $c=g^{X\one}_X$, scalars. By~\eqref{eq:gauge-scalar},
\begin{align*}
\tilde F[X,X,X,X]_{(\one,0,0),(\one,0,0)}&=\frac{ab}{ac}\,F[X,X,X,X]_{(\one,0,0),(\one,0,0)}
=\frac bc\cdot\frac1{d_X},\\
\tilde F[X,\one,X,\one]_{(X,0,0),(X,0,0)}&=\frac{ca}{ba}\,F[X,\one,X,\one]_{(X,0,0),(X,0,0)}=\frac cb,
\end{align*}
using $F[X,X,X,X]_{(\one,0,0),(\one,0,0)}=1/d_X$ and $F[X,\one,X,\one]=1$ in our data. The product
of the two entries is $1/d_X$, of valuation $-2/3$, so they cannot both be integral.
\end{proof}

The object \ob{0011} is the only one whose dimension is divisible by the primes above $13$.

\section{The data files}\label{sec:data}
\subsection{Text format}
All text files are line oriented, with header lines starting with \texttt{\#} that state the field,
the basis and the conventions, list the objects (in the order of Table~\ref{tab:objects}), and give
the fusion multiplicities as lines \texttt{\# N X Y Z n}. The data lines are
\begin{verbatim}
F X Y W T | Z i j | Z' k l | e1 e2 e3 e4 e5 e6
R X Y Z | j i | b0 b1 ... b23
Rinv X Y Z | i j | b0 b1 ... b23
\end{verbatim}
meaning $F[X,Y,W,T]_{(Z,i,j),(Z',k,l)}=\sum_je_j\eta_j$, with $\eta_j=\zeta_{13}^j+\zeta_{13}^{-j}$ and
$\zeta_{13}=\zeta^4$, and $R[X,Y,Z]_{j,i}=\sum_kb_k\zeta^k$ (and the same for $R^{-1}$), with indices from
$0$ and coefficients integers or fractions \texttt{a/b} in lowest terms. The R-symbols lie in $\Kt$, so
only even powers of $\zeta$ occur; they are written in the basis of $\K$ so that the same verifiers read
them. Entries equal to zero are omitted, the lines are sorted by key, and the label sets of each matrix
are determined by the fusion rules. The $\eta_j$ form an integral basis of $\kk$; since
$1=-(\eta_1+\dots+\eta_6)$, the element $1$ is written \texttt{-1 -1 -1 -1 -1 -1}.

\subsection{Binary format}
The complete F-table is stored in a binary format (\texttt{fsym\_bin.py}), in two files of $6.50$ and
$6.45$ GB: a magic string, the length of a JSON header, the header (objects, conventions, and for each
block its key, row and column labels, byte offset and limb counts), and for each block the common
denominator $D$ of its coefficients, the signs, and the magnitudes $|a_kD|$ as little endian $32$-bit
limbs, row major. The coefficients are those of the basis $1,c_1,\dots,c_{11}$ of $\Lf$,
$c_k=\zeta^k+\zeta^{-k}$ (every entry lies in $\kk$), so that the verifiers read the table directly. A
reader returns any block either as exact fractions or reduced modulo a prime at a chosen root, without
reading the rest of the file, and \texttt{fsym\_k.py} prints chosen blocks, or the whole table, in the text
format above:
\begin{verbatim}
python3 code/fields/fsym_k.py archive/F4_4_F_part1.bin \
    archive/F4_4_F_part2.bin 0001,0001,0001,0001
\end{verbatim}

\subsection{Worked examples}
\begin{example}\label{ex:VVVV}
Take $X=Y=W=T=V=\ob{0001}$. Since $V\otimes V=\one\oplus V\oplus\ob{0002}\oplus\ob{0010}\oplus
\ob{1000}$ and each of these has $V$ in its tensor product with $V$ exactly once, the rows and
the columns of $F[V,V,V,V]$ are labelled by $(Z,0,0)$, $Z\in\{\ob{0000},\ob{0001},\ob{0002},
\ob{0010},\ob{1000}\}$, and the matrix is $5\times5$. Its first row is
\begin{verbatim}
F 0001 0001 0001 0001 | 0000 0 0 | 0000 0 0 | 1/2 0 0 0 1/2 0
F 0001 0001 0001 0001 | 0000 0 0 | 0001 0 0 | 1/2 0 0 0 1/2 0
F 0001 0001 0001 0001 | 0000 0 0 | 0002 0 0 | -1 -1 -1 -1 -1 -1
F 0001 0001 0001 0001 | 0000 0 0 | 0010 0 0 | -1 -1 -1 -1 -1 -1
F 0001 0001 0001 0001 | 0000 0 0 | 1000 0 0 | -1 -1 -1 -1 -1 -1
\end{verbatim}
So $F[V,V,V,V]_{(\one,0,0),(\one,0,0)}=\frac12(\eta_1+\eta_5)$, and the other entries of the first row
are $\frac12(\eta_1+\eta_5)$, $1$, $1$, $1$. In $\Lf$, $\eta_j=c_{4j}$ and $c_{20}=-c_6$, so
$\frac12(\eta_1+\eta_5)=\frac12(c_4-c_6)$; since $d_V=2+2c_2+2c_{10}$ and
$\frac12(c_4-c_6)(2+2c_2+2c_{10})=1$, this entry is $1/d_V$, as in Theorem~\ref{thm:main}(iii). At
$\zeta_{13}=e^{2\pi i/13}$ the matrix is approximately
\[
\begin{pmatrix}
0.136945 & 0.136945 & 1 & 1 & 1\\
1 & -0.521601 & 0.785500 & -1.525349 & 4.877810\\
1 & 0.107570 & 1.365971 & -1.087046 & -2.337893\\
1 & -0.208889 & -1.087046 & 1.582221 & -0.936273\\
1 & 0.667993 & -2.337893 & -0.936273 & 2.482147
\end{pmatrix},
\]
with $1/0.136945\approx7.302187=d_V$; the entries outside the first row have integer coefficients.
\end{example}

\begin{example}\label{ex:R}
The R-file contains, for $X=Y=V$,
\begin{verbatim}
R 0001 0001 0000 | 0 0 | 0 0 -1 0 0 0 0 0 0 0 0 0 0 0 0 0 0 0 0 0 0 0 0 0
R 0001 0001 0001 | 0 0 | 0 0 0 0 0 0 0 0 0 0 0 0 0 0 -1 0 0 0 0 0 0 0 0 0
R 0001 0001 0002 | 0 0 | 0 0 1 0 0 0 0 0 0 0 0 0 0 0 0 0 0 0 0 0 0 0 0 0
R 0001 0001 0010 | 0 0 | -1 0 0 0 0 0 0 0 0 0 0 0 0 0 0 0 0 0 0 0 0 0 0 0
R 0001 0001 1000 | 0 0 | 0 0 0 0 0 0 0 0 0 0 0 0 0 0 0 0 0 0 0 0 1 0 0 0
\end{verbatim}
that is $R[V,V,Z]=-\zeta^2,-\zeta^{14},\zeta^2,-1,\zeta^{20}$ for $Z=\one,V,\ob{0002},\ob{0010},\ob{1000}$.
Since $\zeta^{26}=-1$ these are $q^{-24},q^{-12},q^2,-1,-q^{-6}$, the eigenvalues of $c_{V,V}$ of
Section~\ref{sec:cells}, and one checks $\sum_Zd_ZR[V,V,Z]=q^{24}d_V=\theta_Vd_V$. For a multiplicity
space of dimension $2$, the block $R[V,\ob{0011},\ob{0011}]$ consists of four lines
\texttt{R 0001 0011 0011 | j i | ...}, $j,i\in\{0,1\}$; its product with $R[\ob{0011},V,\ob{0011}]$ is
$q^{-24}$ times the identity, as the balancing identity requires.
\end{example}

\section{Availability and reproduction}\label{sec:availability}
\subsection{Where things are}
\begin{itemize}
\item \emph{Programs}: the repository \texttt{kim-em/f4\_4}, \url{https://github.com/kim-em/f4_4}.
Its directory \texttt{code/} has seven subdirectories: \texttt{f4/} (the construction modulo $p$ and the
verifiers), \texttt{fields/} (the normal form, the driver of the pipeline and the reconstruction),
\texttt{checks/} (the exact checks of this paper), \texttt{autoeq/} (Section~\ref{sec:autoeq}),
\texttt{f4-forms/} (Section~\ref{sec:forms}), \texttt{unitary/} (Section~\ref{sec:unitarity}) and
\texttt{metaplectic/} (the comparison for $\cC(\mathfrak{so}(13),2)$ of Section~\ref{sec:limits}). The
repository also contains the logs of all runs (\texttt{logs/}), the deposit (\texttt{deposit/}), a
description of the archive with checksums (\texttt{archive/README.md}) and the source of this paper.
\item \emph{Deposit} (with this paper, about $10$ MB): the files of Table~\ref{tab:deposit}, the logs of
the runs reported here, and a \texttt{README} with SHA-256 checksums.
\item \emph{Archive}: Zenodo, \url{https://doi.org/10.5281/zenodo.24269179}. The complete F-table in the binary
format of Section~\ref{sec:data} (two files \texttt{F4\_4\_F\_part1.bin} and \texttt{F4\_4\_F\_part2.bin},
$6.50$ and $6.45$ GB) and the R-file \texttt{F4\_4\_R\_normalform.txt} (the uncompressed file of the
deposit).
\end{itemize}

\begin{table}[ht]
\centering
\small
\begin{tabular}{lll}
\toprule
file & content & size \\
\midrule
\texttt{F4\_4\_H\_normalform.txt.xz} & F with third leg $V$, over $\kk$; determines the table & 9.0 MB \\
\texttt{F4\_4\_R\_normalform.txt.xz} & R and $R^{-1}$, over $\Kt$ & 0.4 MB \\
\texttt{F4\_4\_invariants.txt} & $d_X$, $\theta_X$, $\tilde S$ exactly (Section~\ref{sec:modular}) & 0.02 MB \\
\bottomrule
\end{tabular}
\medskip
\caption{Data files of the deposit (formats in Section~\ref{sec:data}). The F-symbols with third leg
$V$ determine the whole table (Lemma~\ref{lem:Hdetermines}), which can be regenerated from them by the
pentagon recursion.}
\label{tab:deposit}
\end{table}

\subsection{Programs}
Table~\ref{tab:scripts} lists the programs. They require Python~3 with numpy; the construction also
uses scipy (sparse matrices), and \texttt{exact\_checks.py} uses mpmath. The versions used were
Python~3.14.8, numpy~2.5.3, scipy~1.18.1, sympy~1.14.0 and mpmath~1.4.1, and PARI/GP~2.17.3 for the
\texttt{.gp} scripts. The verifiers
\texttt{verify\_hex\_bin.py}, \texttt{pent\_reduced.py} and \texttt{pent\_exact\_subset.py} use only numpy and
the standard library, and share no code with the programs that produced the data apart from the
reader \texttt{fsym\_bin.py}.

\begin{table}[ht]
\centering
\scriptsize
\renewcommand{\arraystretch}{1.05}
\begin{tabular}{p{0.40\linewidth}p{0.46\linewidth}l}
\toprule
program & role & section \\
\midrule
\multicolumn{3}{p{0.9\linewidth}}{\emph{Construction modulo one prime and one root} (\texttt{code/f4/}, called by \texttt{code/fields/direct\_f4.py})}\\
\texttt{rootsys.py}, \texttt{basic.py} & root data, Weyl and quantum dimensions, Kac--Walton fusion & \ref{sec:category} \\
\texttt{lin.py}, \texttt{qg.py} & linear algebra mod $p$; $\Uq$-modules, highest weight vectors & \ref{sec:modules} \\
\texttt{cells.py}, \texttt{check\_cells.py} & the cells; their relations at every base & \ref{sec:cells} \\
\texttt{pathmodel.py}, \texttt{objects.py} & path spaces and generators; idempotents $p_S$, auxiliary bases & \ref{sec:paths}--\ref{sec:gaugefix} \\
\texttt{hfamily.py}, \texttt{run\_h.py} & F-symbols with third leg $V$ & \ref{sec:H} \\
\texttt{recursion.py}, \texttt{run\_split.py} & the pentagon recursion & \ref{sec:recursion} \\
\texttt{rmatrix.py}, \texttt{run\_r.py} & R-symbols & \ref{sec:Rsym} \\
\multicolumn{3}{p{0.9\linewidth}}{\emph{Normal form, pipeline and reconstruction} (\texttt{code/fields/})}\\
\texttt{f4\_canon.py}, \texttt{torus.py} & the normal form (Section~\ref{sec:nfrule}) modulo $p$ & \ref{sec:nfrule} \\
\texttt{direct\_f4.py}, \texttt{direct\_f4\_all.sh} & the pipeline at one (prime, root); at all of them & \ref{sec:construction} \\
\texttt{compare\_direct.py} & one run of the pipeline against the archived table & \ref{sec:construction} \\
\texttt{f4\_full\_recon.py}, \texttt{direct\_R\_recon.py} & reconstruction of F over $\kk$ and of R over $\Kt$ & \ref{sec:reconstruction} \\
\texttt{finish\_all.sh} & the long verifications; comparison with the archive & \ref{sec:verification} \\
\texttt{fsym\_k.py}, \texttt{fsym\_bin.py} & readers and text output & \ref{sec:data} \\
\texttt{f4\_relabel\_exact.py} & exact relabellings & \ref{sec:autoeq} \\
\multicolumn{3}{p{0.9\linewidth}}{\emph{Verifiers} (\texttt{code/f4/})}\\
\texttt{verify\_hex\_bin.py} & both hexagons and $RR^{-1}=1$, exactly & \ref{sec:verification} \\
\texttt{pent\_reduced.py} & pentagons with fourth object $V$, exactly; invertibility & \ref{sec:pentagon} \\
\texttt{pent\_exact\_subset.py} & pentagon on subsets, exactly & \ref{sec:pentagon} \\
\multicolumn{3}{p{0.9\linewidth}}{\emph{Checks of this paper} (\texttt{code/checks/})}\\
\texttt{exact\_checks.py} & objects, fusion graph, modular data, unit legs, $1/d_X$, exactly & \ref{sec:modular} \\
\texttt{s\_nondegenerate.py} & $\tilde S^2=\mathcal D^2\cdot\id$, exactly & \ref{sec:modular} \\
\texttt{nf\_galois.py}, \texttt{r\_invariants.py} & Galois action on F; invariants of R & \ref{sec:fields} \\
\texttt{pseudo\_unitary.py} & character of $d_X$; $(F^{-1})_{11}=F_{11}$ modulo $p$ & \ref{sec:modularity} \\
\texttt{autoeq\_checks.py}, \texttt{nf\_witness.py} & automorphisms, M\"uger centre, transport of $R[X,X,Z]$; invariants $F_{ab}(F^{-1})_{ba}$ & \ref{sec:autoeq} \\
\multicolumn{3}{p{0.9\linewidth}}{\emph{Other} (\texttt{code/autoeq/}, \texttt{code/unitary/})}\\
\texttt{fusion\_auts.py}, \texttt{stab\_check.py} & fusion ring automorphisms; stabiliser of the normal form & \ref{sec:autoeq} \\
\texttt{braid\_transport\_f4.py} & transported R modulo $p$ & \ref{sec:autoeq} \\
\texttt{f4\_forms.py}, \texttt{f4\_forms\_recon.py} & invariant Hermitian forms modulo $p$; over $\kk$ & \ref{sec:unitarity} \\
\texttt{f4\_forms\_analyse.py}, \texttt{f4\_classes.py} & positivity of the forms; square classes & \ref{sec:unitarity} \\
\multicolumn{3}{p{0.9\linewidth}}{\emph{Forms over the cubic field} (\texttt{code/f4-forms/})}\\
\texttt{relabel\_exact\_nf.py}, \texttt{run\_gauges.sh} & exact gauges for $\sigma_5$ with $\pi$, $\pi^3$, saved & \ref{sec:forms} \\
\texttt{descent\_u.py} & descent data, scalars $u_X$, braiding test & \ref{sec:forms} \\
\texttt{brauer.py}, \texttt{brauer.gp} & $u_X$ in $E$; signs, Hasse invariants, index (PARI/GP) & \ref{sec:forms} \\
\texttt{uvals\_gp.py}, \texttt{endo.gp}, \texttt{kforms.gp} & forms $\mathcal D_e$; twisted forms over $\kk$ (PARI/GP) & \ref{sec:forms} \\
\bottomrule
\end{tabular}
\medskip
\caption{Programs. The comparison for $\cC(\mathfrak{so}(13),2)$ mentioned in
Section~\ref{sec:limits} is in the directory \texttt{code/metaplectic/} of the repository.}
\label{tab:scripts}
\end{table}

\subsection{Commands}
Paths are relative to the top directory of the repository, which contains \texttt{code/} and
\texttt{deposit/}; the archive files are read from the directory named by the environment variable
\texttt{F4\_ARCHIVE} (default \texttt{archive/}), and scratch files are written to \texttt{work/}. Times
and peak memory were measured on a shared $96$-core machine; the times of the long runs are taken from
their logs.

\emph{(a) Verifying the data.} Without the archive:
\begin{verbatim}
python3 code/checks/s_nondegenerate.py deposit/F4_4_invariants.txt
python3 code/checks/r_invariants.py deposit/F4_4_R_normalform.txt.xz
python3 code/checks/autoeq_checks.py deposit/F4_4_R_normalform.txt.xz
python3 code/checks/nf_galois.py deposit/F4_4_H_normalform.txt.xz
\end{verbatim}
These prove $\tilde S^2=\mathcal D^2\cdot\id$ ($2$ s, $30$ MB), the statements about the invariants of R
in Section~\ref{sec:fields} ($12$ s, $61$ MB), the braiding statements of Theorem~\ref{thm:autoeq} that
use only R ($1$ s, $61$ MB) and the minimality of $\kk$ ($13$ s, $0.3$ GB). With the archive:
\begin{verbatim}
A=${F4_ARCHIVE:-archive}; R=$A/F4_4_R_normalform.txt
B="$A/F4_4_F_part1.bin $A/F4_4_F_part2.bin"
python3 code/checks/exact_checks.py deposit/F4_4_R_normalform.txt.xz \
    $B work/inv.txt
python3 code/checks/nf_witness.py $B
python3 code/checks/pseudo_unitary.py $R $B
python3 code/f4/pent_exact_subset.py 0000,0001,1000,0004,2000 $B --jobs 2
python3 code/f4/verify_hex_bin.py $R $B --jobs 2
python3 code/f4/pent_reduced.py $B --jobs 2
\end{verbatim}
\texttt{exact\_checks.py} proves Theorem~\ref{thm:main}(iii)
and the statements of Section~\ref{sec:modular}, and writes the invariants file, identical to the
deposited one ($50$ s, $0.5$ GB); \texttt{nf\_witness.py} gives the monoidal obstructions of
Theorem~\ref{thm:autoeq} ($3$ s, $0.4$ GB); \texttt{pseudo\_unitary.py} checks the character property of the $d_X$ and the
identity of Lemma~\ref{lem:norm} ($30$ s, $0.7$ GB); the fourth command proves the pentagon for the $1{,}585$
five-tuples from the listed objects ($85$ s, $0.9$ GB; with \texttt{maxdim=50} it is the run of
Section~\ref{sec:pentagon}, $21{,}721$ s with $2$ workers). \texttt{verify\_hex\_bin.py} proves both
hexagons and $RR^{-1}=1$ ($8593$ s with $2$ workers) and ends with \texttt{CONCLUSION: both hexagon equations hold exactly in K for all 58740 four-tuples \dots{}}. \texttt{pent\_reduced.py} proves
invertibility and the pentagons with fourth object $V$ and ends with \texttt{CONCLUSION: the pentagon equation holds exactly in K for all 63841 five-tuples with fourth object 0001}
($39{,}892$ s with $2$ workers). The last two write a scratch file of $2.6$ GB per prime.

\emph{(b) One run of the pipeline} at $p=1073741189$, $q=z$:
\begin{verbatim}
python3 code/fields/direct_f4.py 1073741189 1 work/direct --nproc 4
python3 code/fields/compare_direct.py work/direct 1073741189 1
\end{verbatim}
The first command runs steps (1)--(6) of Section~\ref{sec:construction} and writes the residues of the
whole F-table and of R and $R^{-1}$ (at $q=z$ and $q=z^{-1}$) in the normal-form gauge, asserting that
the normal-form rule reproduces the stored choices. Measured: $305$ s with $4$ processes, $1.1$ GB
(cells $6$ s, the family with third leg $V$ $18$ s, the normal form $1$ s, the recursion $182$ s,
export $39$ s, R $58$ s). The second command reduces the archived table at the same root: the residues
agree in all $54{,}881{,}076$ entries ($25$ s, $0.8$ GB).

\emph{(c) Everything.} Both scripts are in \texttt{code/fields} and take a scratch directory as
argument. \texttt{direct\_f4\_all.sh} runs (b) at the $78$ (prime, root) pairs of
Section~\ref{sec:reconstruction} (about $9$ hours with $2$ runs at a time), and
\texttt{f4\_full\_recon.py} and \texttt{direct\_R\_recon.py} reconstruct the F-table over $\kk$ ($279$ s
with $16$ processes) and R over $\Kt$ (\texttt{f4\_full\_recon.py} reads only the block layout of the
archived table). Then \texttt{finish\_all.sh} runs the long verifications of (a) and compares the
reconstructed files with the archived ones: they are identical byte for byte
(\texttt{deposit/logs/finish\_all.log}).

\emph{(d) Forms over the cubic field}, in \texttt{code/f4-forms}:
\begin{verbatim}
./run_gauges.sh
python3 descent_u.py data/h_pi1_a5.pkl data/h_pi3_a5.pkl \
    ../../deposit/F4_4_R_normalform.txt.xz data
python3 brauer.py data/u_values.pkl brauer.gp && gp -q brauer.gp
python3 uvals_gp.py data/u_values.pkl uvals.gp && gp -q endo.gp
\end{verbatim}
The first command computes and proves the two
gauges of Lemma~\ref{lem:twisted} from the F-table ($8.5$ min and $0.65$ GB for each of its two
processes), the second the scalars $u_X$ and the checks of Section~\ref{sec:forms} ($17$ s), the third
the Hasse invariants and indices ($21$ s), the fourth the statements of Proposition~\ref{prop:endo} and the
real forms ($1$ s; \texttt{kforms.gp} takes $23$ s). The logs are in \texttt{deposit/logs/forms\_*.log}.

\section{Limitations and an open question}\label{sec:limits}
\subsection{What the results establish}
The data define a braided fusion category over $\Kt$, with F-symbols over $\kk$ and the fusion rules of
$\CF$ (Theorem~\ref{thm:main}); with its spherical structure of positive dimensions it is a modular tensor
category whose dimensions, twists and $S$-matrix are those of $\CF$ (Proposition~\ref{prop:modular}),
and it has the autoequivalences and braidings of Theorem~\ref{thm:autoeq}.

That this category is $\CF$ rests on the construction from $\Uq$ (in which the braiding on $V\otimes V$ is
taken from the classical eigenvalue formula), together with the agreement of the fusion rules, the
dimensions, the twists and the $S$-matrix with those of $\CF$; we know of no theorem that identifies a
braided fusion category from these invariants. The same programs, applied to
$\cC(\mathfrak{so}(13),2)$ with the spinor as generator, produce F- and
R-symbols that agree exactly, up to an explicit gauge transformation, with the closed formulas of
Ardonne, Finch and Titsworth~\cite{ArdonneFinchTitsworth}\footnote{In arXiv:1608.03762v2 the tables of
F- and R-symbols are correct as printed. For the comparison the following equations are to be read
differently: in the fusion rules of Section~4.1, $\phi_i\otimes\psi_\pm=\psi_+\oplus\psi_-$; in~(12) the
summation index is $g$; the gauge law~(24) for R is $R_{ab}^c\mapsto g_{ba}^cR_{ab}^c/g_{ab}^c$, and the
second hexagon~(17) is~(16) for $(R_{ba}^c)^{-1}$; in~(18) the last factor is an entry of the inverse
F-matrix; in~(51) and~(52) the signs of $S_{\psi_\pm\psi_\pm}$ and $S_{\psi_\pm\psi_\mp}$ are reversed.
The printed R-symbols are the inverses of those in our convention.} and with the solution of SageMath's
F-matrix solver~\cite{SageFMatrix}.

\subsection{Open question}
For which primes $\mathfrak p$ of $\K$ is there a gauge in which all F-symbols of $\CF$ are
$\mathfrak p$-integral? By Proposition~\ref{prop:integrality} the primes above $13$ are not among them.

\section*{Acknowledgements}
The computations, the verification programs and the text of this paper were produced by AI systems
(Claude, by Anthropic) directed by Kim Morrison.

\bibliographystyle{amsplain}
\providecommand{\bysame}{\leavevmode\hbox to3em{\hrulefill}\thinspace}
\providecommand{\MR}{\relax\ifhmode\unskip\space\fi MR }
% \MRhref is called by the amsart/book/proc definition of \MR.
\providecommand{\MRhref}[2]{%
  \href{http://www.ams.org/mathscinet-getitem?mr=#1}{#2}
}
\providecommand{\href}[2]{#2}
\begin{thebibliography}{10}

\bibitem{VerlindeDB}
\emph{{VerlindeDB}: a list of low(ish) rank {V}erlinde categories},
  \url{https://verlindedb.github.io/}, Site source at commit 5aed33f
  (2026-02-11), accessed 2026-11-14; no authors are named on the site.

\bibitem{SageFMatrix}
Willie Aboumrad, \emph{Quantum computing with anyons: an {$F$}-matrix and braid
  calculator}, 2022.

\bibitem{Andersen92}
Henning~Haahr Andersen, \emph{Tensor products of quantized tilting modules},
  Comm. Math. Phys. \textbf{149} (1992), no.~1, 149--159.

\bibitem{AndersenParadowski}
Henning~Haahr Andersen and Jan Paradowski, \emph{Fusion categories arising from
  semisimple {L}ie algebras}, Comm. Math. Phys. \textbf{169} (1995), no.~3,
  563--588.

\bibitem{Alatc}
Eddy Ardonne, \emph{alatc: affine-lie-algebra-tensor-category, a {M}athematica
  package},
  \url{https://github.com/ardonne/affine-lie-algebra-tensor-category}, 2025,
  Last update 2025-04-28 (commit c885a37), accessed 2026-11-14.

\bibitem{ArdonneFinchTitsworth}
Eddy Ardonne, Peter~E. Finch, and Matthew Titsworth, \emph{Classification of
  metaplectic fusion categories}, Symmetry \textbf{13} (2021), no.~11, 2102.

\bibitem{ArdonneSlingerland}
Eddy Ardonne and Joost~K. Slingerland, \emph{Clebsch--{G}ordan and
  6j-coefficients for rank 2 quantum groups}, J. Phys. A \textbf{43} (2010),
  no.~39, 395205.

\bibitem{BakalovKirillov}
Bojko Bakalov and Alexander Kirillov, Jr., \emph{Lectures on tensor categories
  and modular functors}, University Lecture Series, vol.~21, American
  Mathematical Society, Providence, RI, 2001.

\bibitem{Bourbaki}
Nicolas Bourbaki, \emph{Lie groups and {L}ie algebras. {C}hapters 4--6},
  Elements of Mathematics, Springer-Verlag, Berlin, 2002, Translated from the
  French by Andrew Pressley.

\bibitem{DGNO}
Vladimir Drinfeld, Shlomo Gelaki, Dmitri Nikshych, and Victor Ostrik, \emph{On
  braided fusion categories. {I}}, Selecta Math. (N.S.) \textbf{16} (2010),
  no.~1, 1--119, arXiv:0906.0620.

\bibitem{EdieMichell}
Cain Edie-Michell, \emph{Auto-equivalences of the modular tensor categories of
  type {$A$}, {$B$}, {$C$} and {$G$}}, Adv. Math. \textbf{402} (2022), Paper
  No. 108364, With an appendix by T. Gannon; arXiv:2002.03220.

\bibitem{EtingofGelakiDescent}
Pavel Etingof and Shlomo Gelaki, \emph{Descent and forms of tensor categories},
  Int. Math. Res. Not. IMRN (2012), no.~13, 3040--3063, Numbering of results as
  in arXiv:1102.0657v2.

\bibitem{EGNO}
Pavel Etingof, Shlomo Gelaki, Dmitri Nikshych, and Victor Ostrik, \emph{Tensor
  categories}, Mathematical Surveys and Monographs, vol. 205, American
  Mathematical Society, Providence, RI, 2015.

\bibitem{ENO}
Pavel Etingof, Dmitri Nikshych, and Viktor Ostrik, \emph{On fusion categories},
  Ann. of Math. (2) \textbf{162} (2005), no.~2, 581--642.

\bibitem{Finkelberg}
Michael Finkelberg, \emph{An equivalence of fusion categories}, Geom. Funct.
  Anal. \textbf{6} (1996), no.~2, 249--267.

\bibitem{Galindo}
C{\'e}sar Galindo, \emph{On split forms of fusion categories}, 2026,
  arXiv:2608.23743.

\bibitem{Gannon}
Terry Gannon, \emph{The automorphisms of affine fusion rings}, Adv. Math.
  \textbf{165} (2002), no.~2, 165--193, arXiv:math/0002044.

\bibitem{Humphreys}
James~E. Humphreys, \emph{Introduction to {L}ie algebras and representation
  theory}, Graduate Texts in Mathematics, vol.~9, Springer-Verlag, New York,
  1972.

\bibitem{Jantzen}
Jens~Carsten Jantzen, \emph{Lectures on quantum groups}, Graduate Studies in
  Mathematics, vol.~6, American Mathematical Society, Providence, RI, 1996.

\bibitem{KacBook}
Victor~G. Kac, \emph{Infinite-dimensional {L}ie algebras}, third ed., Cambridge
  University Press, Cambridge, 1990.

\bibitem{KacPeterson}
Victor~G. Kac and Dale~H. Peterson, \emph{Infinite-dimensional {L}ie algebras,
  theta functions and modular forms}, Adv. in Math. \textbf{53} (1984), no.~2,
  125--264.

\bibitem{KazhdanLusztigI}
D.~Kazhdan and G.~Lusztig, \emph{Tensor structures arising from affine {L}ie
  algebras. {I}}, J. Amer. Math. Soc. \textbf{6} (1993), no.~4, 905--947.

\bibitem{KazhdanLusztigII}
\bysame, \emph{Tensor structures arising from affine {L}ie algebras. {II}}, J.
  Amer. Math. Soc. \textbf{6} (1993), no.~4, 949--1011.

\bibitem{KazhdanLusztigIII}
\bysame, \emph{Tensor structures arising from affine {L}ie algebras. {III}}, J.
  Amer. Math. Soc. \textbf{7} (1994), no.~2, 335--381.

\bibitem{KazhdanLusztigIV}
\bysame, \emph{Tensor structures arising from affine {L}ie algebras. {IV}}, J.
  Amer. Math. Soc. \textbf{7} (1994), no.~2, 383--453.

\bibitem{MacLane}
Saunders Mac~Lane, \emph{Categories for the working mathematician}, second ed.,
  Graduate Texts in Mathematics, vol.~5, Springer-Verlag, New York, 1998.

\bibitem{TensorCategoriesJl}
Fabian M{\"a}urer, \emph{{TensorCategories.jl}},
  \url{https://github.com/TensorCategories/TensorCategories.jl}, 2026, Version
  0.8.0 (commit c820a48, 2026-11-09), Zenodo doi:10.5281/zenodo.24081897;
  accessed 2026-11-14.

\bibitem{MaurerThiel}
Fabian M{\"a}urer and Ulrich Thiel, \emph{Computing the center of a fusion
  category}, 2024.

\bibitem{MaurerThielVercleyen}
Fabian M{\"a}urer, Ulrich Thiel, and Gert Vercleyen, \emph{F-symbols and
  {R}-symbols for the {D}rinfeld center of the {H}aagerup subfactor}, 2026.

\bibitem{MorrisonSnyderThurston}
Kim Morrison, Noah Snyder, and Dylan~P. Thurston, \emph{Towards the quantum
  exceptional series}, 2024, arXiv:2402.03637.

\bibitem{Rowell}
Eric~C. Rowell, \emph{From quantum groups to unitary modular tensor
  categories}, Representations of algebraic groups, quantum groups, and {L}ie
  algebras, Contemp. Math., vol. 413, American Mathematical Society,
  Providence, RI, 2006, pp.~215--230.

\bibitem{Stasheff}
James~Dillon Stasheff, \emph{Homotopy associativity of {$H$}-spaces. {I}},
  Trans. Amer. Math. Soc. \textbf{108} (1963), no.~2, 275--292.

\bibitem{VercleyenThesis}
Gert Vercleyen, \emph{On low-rank multiplicity-free fusion categories}, Ph.D.
  thesis, 2024.

\bibitem{AnyonWiki}
Gert Vercleyen and Joost Slingerland, \emph{{AnyonWiki}},
  \url{https://anyonwiki.github.io/}, 2023, Version 1.0.2, Zenodo,
  doi:10.5281/zenodo.15863113; site source at commit 25147de (2026-09-03),
  accessed 2026-11-14.

\bibitem{MCA}
Joachim von~zur Gathen and J{\"u}rgen Gerhard, \emph{Modern computer algebra},
  third ed., Cambridge University Press, Cambridge, 2013.

\bibitem{Walton90a}
Mark~A. Walton, \emph{Algorithm for {WZW} fusion rules: a proof}, Phys. Lett. B
  \textbf{241} (1990), no.~3, 365--368.

\bibitem{Walton90b}
\bysame, \emph{Fusion rules in {W}ess--{Z}umino--{W}itten models}, Nuclear
  Phys. B \textbf{340} (1990), no.~2-3, 777--790.

\bibitem{WangGuyDavenport}
Paul~S. Wang, M.~J.~T. Guy, and J.~H. Davenport, \emph{P-adic reconstruction of
  rational numbers}, ACM SIGSAM Bull. \textbf{16} (1982), no.~2, 2--3.

\bibitem{Wenzl}
Hans Wenzl, \emph{{$C^*$} tensor categories from quantum groups}, J. Amer.
  Math. Soc. \textbf{11} (1998), no.~2, 261--282.

\bibitem{Xu}
Feng Xu, \emph{Standard {$\lambda$}-lattices from quantum groups}, Invent.
  Math. \textbf{134} (1998), no.~3, 455--487.

\end{thebibliography}

Document prepared in \href{https://verso.lean-lang.org/}{Verso}.

\end{document}
